\newMonth{Aprilis}
\setrunningtitles{Aprilis}{April}

% ---- martyrology/mart04/mart0401.htm
\needspace{10\baselineskip}
\begin{paracol}{2}
\selectlanguage{latin}
\begin{center}{\color{gregoriocolor} Kaléndis Aprílis. 
 Luna\dots\ }\end{center}
\switchcolumn
\selectlanguage{english}
\begin{center}{\color{gregoriocolor} The 1\textsuperscript{st} Day of April. 
 The\dots\ Day of the Moon.}\end{center}
\end{paracol}

\noindent\begin{tabularx}{\linewidth}{*{19}{>{\centering\arraybackslash}X}}
 \textcolor{gregoriocolor}{a} & \textcolor{gregoriocolor}{b} & \textcolor{gregoriocolor}{c} & \textcolor{gregoriocolor}{d} & \textcolor{gregoriocolor}{e} & \textcolor{gregoriocolor}{f} & \textcolor{gregoriocolor}{g} & \textcolor{gregoriocolor}{h} & \textcolor{gregoriocolor}{i} & \textcolor{gregoriocolor}{k} & \textcolor{gregoriocolor}{l} & \textcolor{gregoriocolor}{m} & \textcolor{gregoriocolor}{n} & \textcolor{gregoriocolor}{p} & \textcolor{gregoriocolor}{q} & \textcolor{gregoriocolor}{r} & \textcolor{gregoriocolor}{s} & \textcolor{gregoriocolor}{t} & \textcolor{gregoriocolor}{u} \\
 3 & 4 & 5 & 6 & 7 & 8 & 9 & 10 & 11 & 12 & 13 & 14 & 15 & 16 & 17 & 18 & 19 & 20 & 21 \\
\end{tabularx}
\vspace{0.5\baselineskip}
\noindent\begin{tabularx}{\linewidth}{*{12}{>{\centering\arraybackslash}X}}
 \textcolor{gregoriocolor}{A} & \textcolor{gregoriocolor}{B} & \textcolor{gregoriocolor}{C} & \textcolor{gregoriocolor}{D} & \textcolor{gregoriocolor}{E} & F & \textcolor{gregoriocolor}{F} & \textcolor{gregoriocolor}{G} & \textcolor{gregoriocolor}{H} & \textcolor{gregoriocolor}{M} & \textcolor{gregoriocolor}{N} & \textcolor{gregoriocolor}{P} \\
 22 & 23 & 24 & 25 & 26 & 27 & 27 & 28 & 29 & 30 & 1 & 2 \\
\end{tabularx}

\begin{paracol}{2}
\selectlanguage{latin}
\lettrine[lines=2]{R}{omæ} pássio sanctæ Theodóræ, soróris 
 illustríssimi Mártyris Hermétis, quæ, sub Hadriáno Imperatóre, ab Aureliáno 
 Júdice affécta martyrio, sepúlta est juxta fratrem suum, via Salária, non 
 longe ab Urbe.
\switchcolumn
\selectlanguage{english}
\lettrine[lines=2]{A}{t} Rome, the passion of St. Theodora, sister of the illustrious martyr 
 Hermes. She underwent martyrdom in the time of Emperor Adrian, under 
 the judge Aurelian, and was buried at the side of her brother, on the 
 Salarian Way, a short distance from the city.
\switchcolumn*
\selectlanguage{latin}
Eódem die sancti Venántii, Epíscopi et Mártyris.
\switchcolumn
\selectlanguage{english}
The same day, St. Venantius, bishop and martyr.
\switchcolumn*
\selectlanguage{latin}
In Ægypto sanctórum Mártyrum Victóris et 
 Stéphani.
\switchcolumn
\selectlanguage{english}
In Egypt, the holy martyrs Victor and Stephen.
\switchcolumn*
\selectlanguage{latin}
In Arménia sanctórum Mártyrum Quinctiáni et Irenǽi.
\switchcolumn
\selectlanguage{english}
In Armenia, the holy martyrs Quinctian and Irenæus.
\switchcolumn*
\selectlanguage{latin}
Constantinópoli sancti Macárii Confessóris, qui, sub Leóne Imperatóre, pro 
 assertióne sanctárum Imáginum, in exsílio vitam finívit.
\switchcolumn
\selectlanguage{english}
At Constantinople, under Emperor Leo, St. Macarius, confessor, who ended his 
 life in exile for defending the veneration of sacred images.
\switchcolumn*
\selectlanguage{latin}
Ardpatrícii, in Momónia Hibérniæ província, 
 sancti Celsi Epíscopi, qui beátum Malachíam in Episcopátu præcéssit.
\switchcolumn
\selectlanguage{english}
At Ard-Patrick in Munster, a province of Ireland, Bishop St. Celsus, who 
 preceded blessed Malachy in that bishopric.
\switchcolumn*
\selectlanguage{latin}
Gratianópoli, in Gállia, sancti Hugónis Epíscopi, qui multis annis in 
 solitúdine vitam exégit, et miraculórum glória clarus migrávit ad Dóminum.
\switchcolumn
\selectlanguage{english}
At Grenoble in France, Bishop St. Hugh, who spent many years of his life in 
 solitude, and departed for heaven with a great reputation for miracles.
\switchcolumn*
\selectlanguage{latin}
Apud Ambiánum, in Gállia, sancti Waleríci 
 Abbátis, cujus sepúlcrum crebris miráculis illustrátur.
\switchcolumn
\selectlanguage{english}
At Amiens in France, Abbot St. Valery, whose tomb is well known for its 
 frequent miracles.
\switchcolumn*
\selectlanguage{latin}
\end{paracol}

% \continueMonth{Aprilis}

% ---- martyrology/mart04/mart0402.htm
\needspace{10\baselineskip}
\begin{paracol}{2}
\selectlanguage{latin}
\begin{center}{\color{gregoriocolor} Quarto Nonas Aprílis. 
 Luna\dots\ }\end{center}
\switchcolumn
\selectlanguage{english}
\begin{center}{\color{gregoriocolor} The 2\textsuperscript{nd} Day of April. 
 The\dots\ Day of the Moon.}\end{center}
\end{paracol}

\noindent\begin{tabularx}{\linewidth}{*{19}{>{\centering\arraybackslash}X}}
 \textcolor{gregoriocolor}{a} & \textcolor{gregoriocolor}{b} & \textcolor{gregoriocolor}{c} & \textcolor{gregoriocolor}{d} & \textcolor{gregoriocolor}{e} & \textcolor{gregoriocolor}{f} & \textcolor{gregoriocolor}{g} & \textcolor{gregoriocolor}{h} & \textcolor{gregoriocolor}{i} & \textcolor{gregoriocolor}{k} & \textcolor{gregoriocolor}{l} & \textcolor{gregoriocolor}{m} & \textcolor{gregoriocolor}{n} & \textcolor{gregoriocolor}{p} & \textcolor{gregoriocolor}{q} & \textcolor{gregoriocolor}{r} & \textcolor{gregoriocolor}{s} & \textcolor{gregoriocolor}{t} & \textcolor{gregoriocolor}{u} \\
 4 & 5 & 6 & 7 & 8 & 9 & 10 & 11 & 12 & 13 & 14 & 15 & 16 & 17 & 18 & 19 & 20 & 21 & 22 \\
\end{tabularx}
\vspace{0.5\baselineskip}
\noindent\begin{tabularx}{\linewidth}{*{12}{>{\centering\arraybackslash}X}}
 \textcolor{gregoriocolor}{A} & \textcolor{gregoriocolor}{B} & \textcolor{gregoriocolor}{C} & \textcolor{gregoriocolor}{D} & \textcolor{gregoriocolor}{E} & F & \textcolor{gregoriocolor}{F} & \textcolor{gregoriocolor}{G} & \textcolor{gregoriocolor}{H} & \textcolor{gregoriocolor}{M} & \textcolor{gregoriocolor}{N} & \textcolor{gregoriocolor}{P} \\
 23 & 24 & 25 & 26 & 27 & 28 & 28 & 29 & 30 & 1 & 2 & 3 \\
\end{tabularx}

\begin{paracol}{2}
\selectlanguage{latin}
\lettrine[lines=2]{T}{urónis,} in Gállia, sancti Francísci de Paula Confessóris, qui Ordinis 
 Minimórum Institútor éxstitit; atque, virtútibus et miráculis clarus, a 
 Leóne Papa Décimo in Sanctórum númerum est adscríptus.
\switchcolumn
\selectlanguage{english}
\lettrine[lines=2]{A}{t} Tours in France, St. Francis of Paula, founder of the Order of Minims. 
 Because he was renowned for virtues and miracles, he was inscribed among the 
 saints by Pope Leo X.
\switchcolumn*
\selectlanguage{latin}
Cæsaréæ, in Palæstína, natális sancti Apphiáni 
 Mártyris, qui, ante sanctum Ædésium Mártyrem, fratrem suum, in persecutióne 
 Galérii Maximiáni, cum Prǽsidem Urbánum idólis immolántem arguísset, sæve 
 dilaniátus est, et, pédibus lino in óleum intíncto eóque incénso obvolútis, 
 acerbíssime cruciátus, ac demum in mare demérsus; atque ita, tránsiens per 
 ignem et aquam, edúctus est in refrigérium.
\switchcolumn
\selectlanguage{english}
At Caesarea in Palestine, during the persecution of Galerius Maximian, the 
 birthday of the martyr St. Amphian. He reproved the governor Urban for 
 sacrificing to idols, so his body was cruelly cut in shreds, his feet 
 wrapped in oil-soaked cloths, and set on fire. After these painful 
 torments, he was cast into the sea. Thus through fire and water, he 
 reached his everlasting repose.
\switchcolumn*
\selectlanguage{latin}
Ibídem pássio sanctæ Theodósiæ, Vírginis Tyriæ, 
 quæ, in eádem persecutióne, cum sanctos Confessóres ante tribúnal stantes 
 públice salutásset, atque rogásset eos, ut, cum ad Dóminum perveníssent, sui 
 recordaréntur, a milítibus tenta et ad Urbánum Prǽsidem ducta est; eóque 
 jubénte, latéribus et mammis ad interióra usque dilaniátis, in mare tandem 
 projícitur.
\switchcolumn
\selectlanguage{english}
In the same city, the passion of St. Theodosia, a virgin of Tyre. In 
 the same persecution, she publicly spoke to the holy confessors as they 
 stood before the tribunal, and begged of them to remember her when they 
 should be with God. She was arrested and led to the governor Urban, at 
 whose order her sides and breasts were deeply lacerated, and she was thrown 
 into the sea.
\switchcolumn*
\selectlanguage{latin}
Apud Língonas, in Gállia, sancti Urbáni 
 Epíscopi.
\switchcolumn
\selectlanguage{english}
At Langres in France, Bishop St. Urban.
\switchcolumn*
\selectlanguage{latin}
Apud Comum sancti Abúndii, Epíscopi et Confessóris.
\switchcolumn
\selectlanguage{english}
At Como, St. Abundius, bishop and confessor.
\switchcolumn*
\selectlanguage{latin}
Cápuæ sancti Victóris Epíscopi, eruditióne et 
 sanctitáte conspícui.
\switchcolumn
\selectlanguage{english}
At Capua, Bishop St. Victor, well known for his sanctity and learning.
\switchcolumn*
\selectlanguage{latin}
Lugdúni, in Gállia, sancti Nicétii, ejúsdem urbis Epíscopi, vita et 
 miráculis clari.
\switchcolumn
\selectlanguage{english}
At Lyons in France, St. Nicetus, bishop of that city, renowned for his life 
 and miracles.
\switchcolumn*
\selectlanguage{latin}
In Palæstína deposítio sanctæ Maríæ Ægyptíacæ, 
 quæ peccátrix appellátur.
\switchcolumn
\selectlanguage{english}
In Palestine, the death of St. Mary of Egypt, called the Sinner.
\switchcolumn*
\selectlanguage{latin}
\end{paracol}


% ---- martyrology/mart04/mart0403.htm
\needspace{10\baselineskip}
\begin{paracol}{2}
\selectlanguage{latin}
\begin{center}{\color{gregoriocolor} Tértio Nonas Aprílis. 
 Luna\dots\ }\end{center}
\switchcolumn
\selectlanguage{english}
\begin{center}{\color{gregoriocolor} The 3\textsuperscript{rd} Day of April. 
 The\dots\ Day of the Moon.}\end{center}
\end{paracol}

\noindent\begin{tabularx}{\linewidth}{*{19}{>{\centering\arraybackslash}X}}
 \textcolor{gregoriocolor}{a} & \textcolor{gregoriocolor}{b} & \textcolor{gregoriocolor}{c} & \textcolor{gregoriocolor}{d} & \textcolor{gregoriocolor}{e} & \textcolor{gregoriocolor}{f} & \textcolor{gregoriocolor}{g} & \textcolor{gregoriocolor}{h} & \textcolor{gregoriocolor}{i} & \textcolor{gregoriocolor}{k} & \textcolor{gregoriocolor}{l} & \textcolor{gregoriocolor}{m} & \textcolor{gregoriocolor}{n} & \textcolor{gregoriocolor}{p} & \textcolor{gregoriocolor}{q} & \textcolor{gregoriocolor}{r} & \textcolor{gregoriocolor}{s} & \textcolor{gregoriocolor}{t} & \textcolor{gregoriocolor}{u} \\
 5 & 6 & 7 & 8 & 9 & 10 & 11 & 12 & 13 & 14 & 15 & 16 & 17 & 18 & 19 & 20 & 21 & 22 & 23 \\
\end{tabularx}
\vspace{0.5\baselineskip}
\noindent\begin{tabularx}{\linewidth}{*{12}{>{\centering\arraybackslash}X}}
 \textcolor{gregoriocolor}{A} & \textcolor{gregoriocolor}{B} & \textcolor{gregoriocolor}{C} & \textcolor{gregoriocolor}{D} & \textcolor{gregoriocolor}{E} & F & \textcolor{gregoriocolor}{F} & \textcolor{gregoriocolor}{G} & \textcolor{gregoriocolor}{H} & \textcolor{gregoriocolor}{M} & \textcolor{gregoriocolor}{N} & \textcolor{gregoriocolor}{P} \\
 24 & 25 & 26 & 27 & 28 & 29 & 29 & 30 & 1 & 2 & 3 & 4 \\
\end{tabularx}

\begin{paracol}{2}
\selectlanguage{latin}
\lettrine[lines=2]{R}{omæ} natális beáti Xysti Primi, Papæ et 
 Mártyris; qui, tempóribus Hadriáni Imperatóris, summa cum laude rexit 
 Ecclésiam, ac demum, sub Antoníno Pio, ut sibi Christum lucrifáceret, 
 libénter mortem sustínuit temporálem.
\switchcolumn
\selectlanguage{english}
\lettrine[lines=2]{A}{t} Rome, the birthday of blessed Pope Sixtus the First, martyr, who ruled 
 the Church with distinction during the reign of Emperor Hadrian, and finally 
 in the reign of Antoninus Pius he gladly accepted temporal death in order to 
 gain Christ for himself.
\switchcolumn*
\selectlanguage{latin}
Tauroménii, in Sicília, sancti Pancrátii Epíscopi, qui Christi Evangélium, 
 quod a sancto Petro Apóstolo illuc missus prædicáverat, 
 martyrii sánguine consignávit.
\switchcolumn
\selectlanguage{english}
At Taormina in Sicily, Bishop St. Pancras, who sealed with a martyr's blood 
 the Gospel of Christ that the apostle St. Peter had sent him there to 
 preach.
\switchcolumn*
\selectlanguage{latin}
Tomis, in Scythia, natális sanctórum Mártyrum Evágrii et Benígni.
\switchcolumn
\selectlanguage{english}
At Tomis in Scythia, the birthday of the holy martyrs Evagrius and Benignus.
\switchcolumn*
\selectlanguage{latin}
Tyri, in Phœnícia, sancti Vulpiáni Mártyris, 
 qui, in persecutióne Maximiáni Galérii, cum áspide et cane insútus cúleo, in 
 mare demérsus fuit.
\switchcolumn
\selectlanguage{english}
At Tyre, the martyr St. Vulpian, who was sewn up in a sack with a serpent 
 and a dog and drowned in the sea, during the persecution of Maximian 
 Galerius.
\switchcolumn*
\selectlanguage{latin}
Thessalonícæ 
 pássio sanctárum Vírginum Agapis et Chióniæ, Diocletiáno Imperatóre, sub quo 
 et sancta Virgo Iréne, eárum soror, póstmodum passúra erat. Ambæ vero, 
 cum Christum negáre nollent, primum in cárcere macerátæ sunt, póstea in 
 ignem missæ, sed a flammis intáctæ, ibi, oratióne ad Dóminum fusa, ánimas 
 reddidérunt.
\switchcolumn
\selectlanguage{english}
At Thessalonica, the martyrdom of the holy virgins Agape and Chionia, under 
 Emperor Diocletian. Because they would not deny Christ, they were 
 first detained in prison, then cast into the fire where, untouched by the 
 flames, they gave up their souls to their Creator while praying. Their 
 sister Irene had been imprisoned with them, but was to die later.
\switchcolumn*
\selectlanguage{latin}
In monastério Medícii, in Bithynia, deposítio sancti Nicétæ 
 Abbátis, qui ob cultum sanctárum Imáginum, sub Leóne Arméno, multa passus 
 est, ac tandem, juxta Constantinópolim, Conféssor quiévit in pace.
\switchcolumn
\selectlanguage{english}
In the monastery of Medicion in Bithynia, Abbot St. Nicetas, who suffered a 
 great deal for the veneration of sacred images in the time of Leo the 
 Armenian, and then died in peace as a confessor near Constantinople.
\switchcolumn*
\selectlanguage{latin}
In Anglia sancti Richárdii, Epíscopi Cicestrénsis, sanctitáte et miraculórum 
 glória conspícui.
\switchcolumn
\selectlanguage{english}
In England, St. Richard, bishop of Chichester, celebrated for his sanctity 
 and glorious miracles.
\switchcolumn*
\selectlanguage{latin}
Eboríaci, in território Meldénsi, sanctæ 
 Burgundofáræ, étiam Faræ nómine appellátæ, Abbatíssæ et Vírginis.
\switchcolumn
\selectlanguage{english}
At Faremoutiers, in the district of Meaux, St. Burgundofara, also known as 
 St. Fara, abbess and virgin.
\switchcolumn*
\selectlanguage{latin}
\end{paracol}


% ---- martyrology/mart04/mart0404.htm
\needspace{10\baselineskip}
\begin{paracol}{2}
\selectlanguage{latin}
\begin{center}{\color{gregoriocolor} Prídie Nonas Aprílis. 
 Luna\dots\ }\end{center}
\switchcolumn
\selectlanguage{english}
\begin{center}{\color{gregoriocolor} The 4\textsuperscript{th} Day of April. 
 The\dots\ Day of the Moon.}\end{center}
\end{paracol}

\noindent\begin{tabularx}{\linewidth}{*{19}{>{\centering\arraybackslash}X}}
 \textcolor{gregoriocolor}{a} & \textcolor{gregoriocolor}{b} & \textcolor{gregoriocolor}{c} & \textcolor{gregoriocolor}{d} & \textcolor{gregoriocolor}{e} & \textcolor{gregoriocolor}{f} & \textcolor{gregoriocolor}{g} & \textcolor{gregoriocolor}{h} & \textcolor{gregoriocolor}{i} & \textcolor{gregoriocolor}{k} & \textcolor{gregoriocolor}{l} & \textcolor{gregoriocolor}{m} & \textcolor{gregoriocolor}{n} & \textcolor{gregoriocolor}{p} & \textcolor{gregoriocolor}{q} & \textcolor{gregoriocolor}{r} & \textcolor{gregoriocolor}{s} & \textcolor{gregoriocolor}{t} & \textcolor{gregoriocolor}{u} \\
 6 & 7 & 8 & 9 & 10 & 11 & 12 & 13 & 14 & 15 & 16 & 17 & 18 & 19 & 20 & 21 & 22 & 23 & 24 \\
\end{tabularx}
\vspace{0.5\baselineskip}
\noindent\begin{tabularx}{\linewidth}{*{12}{>{\centering\arraybackslash}X}}
 \textcolor{gregoriocolor}{A} & \textcolor{gregoriocolor}{B} & \textcolor{gregoriocolor}{C} & \textcolor{gregoriocolor}{D} & \textcolor{gregoriocolor}{E} & F & \textcolor{gregoriocolor}{F} & \textcolor{gregoriocolor}{G} & \textcolor{gregoriocolor}{H} & \textcolor{gregoriocolor}{M} & \textcolor{gregoriocolor}{N} & \textcolor{gregoriocolor}{P} \\
 25 & 26 & 27 & 28 & 29 & 1 & 30 & 1 & 2 & 3 & 4 & 5 \\
\end{tabularx}

\begin{paracol}{2}
\selectlanguage{latin}
\lettrine[lines=2]{H}{íspali,} in Hispánia, sancti Isidóri Epíscopi, 
 Confessóris et Ecclésiæ Doctóris, sanctitáte et doctrína conspícui; qui zelo 
 cathólicæ fidei et ecclesiásticæ observántia disciplinæ Hispánias 
 illustrávit.
\switchcolumn
\selectlanguage{english}
\lettrine[lines=2]{A}{t} Seville in Spain, St. Isidore, bishop, confessor, and doctor of the 
 Church. He was conspicuous for sanctity and learning, and had 
 brightened all Spain by his zeal for the Catholic faith and his observance 
 of Church discipline.
\switchcolumn*
\selectlanguage{latin}
Medioláni deposítio sancti Ambrósii Epíscopi, Confessóris et Ecclésiæ 
 Doctóris; cujus stúdio, inter cétera doctrínæ et miraculórum insígnia, 
 témpore Ariánæ perfídiæ, tota fere Itália ad cathólicam fidem convérsa est. 
 Ipsíus tamen festívitas séptimo Idus Decémbris potíssimum recólitur, quo die 
 Epíscopus Mediolanénsis ordinátus est.
\switchcolumn
\selectlanguage{english}
At Milan, the death of St. Ambrose, bishop and confessor, doctor of the 
 Church. By his zeal, besides other monuments to his learning and 
 miracles, almost all Italy returned to the Catholic faith at the time of the 
 Arian heresy. His feast is properly kept on the seventh of December, 
 on which day he became Bishop of Milan.
\switchcolumn*
\selectlanguage{latin}
Thessalonícæ sanctórum Mártyrum Agathópodis 
 Diáconi, et Theodúli Lectóris, qui, sub Maximiáno Imperatóre et Faustíno Prǽside, ob Christiánæ fidei confessiónem, in mare, 
 alligáto ad collum saxo, 
 demérsi sunt.
\switchcolumn
\selectlanguage{english}
At Thessalonica, in the time of Emperor Maximian and the governor Faustinus, 
 the holy martyrs Agathopodes, a deacon, and Theodulus, a lector, who, for 
 the confession of the Catholic faith, had stones tied to their necks and 
 were drowned in the sea.
\switchcolumn*
\selectlanguage{latin}
Constantinópoli sancti Platónis Mónachi, qui plúribus annis advérsus hæréticos, 
 sanctárum Imáginum effractóres, invícto ánimo decertávit.
\switchcolumn
\selectlanguage{english}
At Constantinople, the monk St. Plato. For many years he combated with 
 dauntless courage the heretics bent on destroying sacred images.
\switchcolumn*
\selectlanguage{latin}
In Palæstína sancti Zósimi Anachorétæ, qui 
 funus sanctæ Maríæ Ægyptíacæ curávit.
\switchcolumn
\selectlanguage{english}
In Palestine, the anchoret St. Zosimus, who took care of the funeral 
 of St. Mary of Egypt.
\switchcolumn*
\selectlanguage{latin}
Panórmi sancti Benedícti a sancto Philadélpho, ob córporis nigrédinem 
 cognoménto Nigri, ex Ordine Minórum, Confessóris; qui, signis et virtútibus 
 clarus, in Dómino quiévit, et a Pio Séptimo, Pontífice Máximo, in Sanctórum 
 númerum relátus est.
\switchcolumn
\selectlanguage{english}
At Palermo, St. Benedict of St. Philadelphus, called the Black because of 
 the darkness of his body, a confessor of the Order of Friars Minor. 
 After becoming outstanding for signs and virtues, he went to rest in the 
 Lord, and was enrolled among the saints by Pope Pius VII.
\switchcolumn*
\selectlanguage{latin}
\end{paracol}


% ---- martyrology/mart04/mart0405.htm
\needspace{10\baselineskip}
\begin{paracol}{2}
\selectlanguage{latin}
\begin{center}{\color{gregoriocolor} Nonis Aprílis. 
 Luna\dots\ }\end{center}
\switchcolumn
\selectlanguage{english}
\begin{center}{\color{gregoriocolor} The 5\textsuperscript{th} Day of April. 
 The\dots\ Day of the Moon.}\end{center}
\end{paracol}

\noindent\begin{tabularx}{\linewidth}{*{19}{>{\centering\arraybackslash}X}}
 \textcolor{gregoriocolor}{a} & \textcolor{gregoriocolor}{b} & \textcolor{gregoriocolor}{c} & \textcolor{gregoriocolor}{d} & \textcolor{gregoriocolor}{e} & \textcolor{gregoriocolor}{f} & \textcolor{gregoriocolor}{g} & \textcolor{gregoriocolor}{h} & \textcolor{gregoriocolor}{i} & \textcolor{gregoriocolor}{k} & \textcolor{gregoriocolor}{l} & \textcolor{gregoriocolor}{m} & \textcolor{gregoriocolor}{n} & \textcolor{gregoriocolor}{p} & \textcolor{gregoriocolor}{q} & \textcolor{gregoriocolor}{r} & \textcolor{gregoriocolor}{s} & \textcolor{gregoriocolor}{t} & \textcolor{gregoriocolor}{u} \\
 7 & 8 & 9 & 10 & 11 & 12 & 13 & 14 & 15 & 16 & 17 & 18 & 19 & 20 & 21 & 22 & 23 & 24 & 25 \\
\end{tabularx}
\vspace{0.5\baselineskip}
\noindent\begin{tabularx}{\linewidth}{*{12}{>{\centering\arraybackslash}X}}
 \textcolor{gregoriocolor}{A} & \textcolor{gregoriocolor}{B} & \textcolor{gregoriocolor}{C} & \textcolor{gregoriocolor}{D} & \textcolor{gregoriocolor}{E} & F & \textcolor{gregoriocolor}{F} & \textcolor{gregoriocolor}{G} & \textcolor{gregoriocolor}{H} & \textcolor{gregoriocolor}{M} & \textcolor{gregoriocolor}{N} & \textcolor{gregoriocolor}{P} \\
 26 & 27 & 28 & 29 & 1 & 2 & 1 & 2 & 3 & 4 & 5 & 6 \\
\end{tabularx}

\begin{paracol}{2}
\selectlanguage{latin}
\lettrine[lines=2]{V}{enétiæ,} in Británnia minóre, sancti Vincéntii, 
 cognoménto Ferrérii, ex Ordine Prædicatórum, Confessóris, qui, potens ópere 
 et sermóne, multa míllia infidélium convértit ad Christum.
\switchcolumn
\selectlanguage{english}
\lettrine[lines=2]{A}{t} Vannes in Brittany, St. Vincent Ferrer, of the Order of Preachers, and 
 confessor. He was mighty in word and deed, and converted many 
 thousands of infidels to Christ.
\switchcolumn*
\selectlanguage{latin}
In Africa pássio sanctórum Mártyrum, qui, in persecutióne Regis Ariáni 
 Genseríci, in die Paschæ, in Ecclésia cæsi sunt, 
 quorum Lector, dum in púlpito « ALLELUJA » cantáret, sagítta in gútture 
 transfíxus est.
\switchcolumn
\selectlanguage{english}
In Africa, during the persecution of the Arian king Genseric, the holy 
 martyrs who were murdered in the church on Easter day. The lector, 
 while singing ``Alleluia'` at the lectern, was pierced through the throat by 
 an arrow.
\switchcolumn*
\selectlanguage{latin}
Eódem die sancti Zenónis Mártyris, qui, pice íllitus et in ignem conjéctus, 
 et hasta in médio rogo vulnerátus, martyrio coronátus est.
\switchcolumn
\selectlanguage{english}
The same day, the martyr St. Zeno, who was covered with pitch, cast into the 
 fire, and wounded by the thrust of a spear, thus gaining the crown of 
 martyrdom.
\switchcolumn*
\selectlanguage{latin}
In Lesbo ínsula pássio sanctárum quinque vírginum, quæ 
 gládio martyrium consummárunt.
\switchcolumn
\selectlanguage{english}
On the island of Lesbos, the martyrdom of five holy virgins, who were slain 
 by the sword.
\switchcolumn*
\selectlanguage{latin}
Thessalonícæ sanctæ Irénes Vírginis, quæ, cum 
 sacros Libros contra Diocletiáni edíctum occultásset, ideo, post tolerántiam cárceris, sagítta percússa est et igne cremáta, jussu Dulcétii Præsidis; sub 
 quo et soróres ejus simul Agape et Chiónia passæ ántea fúerant.
\switchcolumn
\selectlanguage{english}
At Thessalonica, the virgin St. Irene, who was imprisoned for hiding the 
 sacred books, contrary to the order of Diocletian. She was pierced 
 with an arrow, then burned to death by order of the governor Dulcetius, 
 under whom her sisters Agape and Chionia had previously suffered.
\switchcolumn*
\selectlanguage{latin}
Palmæ, in ínsula Majórica, sanctæ Catharínæ 
 Thomas, Vírginis, Canoníssæ Reguláris, ex Ordine sancti Augustíni, quam Pius 
 Papa Undécimus in sanctárum Vírginum númerum rétulit.
\switchcolumn
\selectlanguage{english}
In the monastery at Palma, in the diocese of Majorca, the birthday of St. 
 Catalina Tomas, Canoness Regular of the Order of St. Augustine, whom Pope 
 Pius XI, in the fiftieth year of his priesthood, placed among the number of 
 virgin saints.
\switchcolumn*
\selectlanguage{latin}
\end{paracol}


% ---- martyrology/mart04/mart0406.htm
\needspace{10\baselineskip}
\begin{paracol}{2}
\selectlanguage{latin}
\begin{center}{\color{gregoriocolor} Octávo Idus Aprílis. 
 Luna\dots\ }\end{center}
\switchcolumn
\selectlanguage{english}
\begin{center}{\color{gregoriocolor} The 6\textsuperscript{th} Day of April. 
 The\dots\ Day of the Moon.}\end{center}
\end{paracol}

\noindent\begin{tabularx}{\linewidth}{*{19}{>{\centering\arraybackslash}X}}
 \textcolor{gregoriocolor}{a} & \textcolor{gregoriocolor}{b} & \textcolor{gregoriocolor}{c} & \textcolor{gregoriocolor}{d} & \textcolor{gregoriocolor}{e} & \textcolor{gregoriocolor}{f} & \textcolor{gregoriocolor}{g} & \textcolor{gregoriocolor}{h} & \textcolor{gregoriocolor}{i} & \textcolor{gregoriocolor}{k} & \textcolor{gregoriocolor}{l} & \textcolor{gregoriocolor}{m} & \textcolor{gregoriocolor}{n} & \textcolor{gregoriocolor}{p} & \textcolor{gregoriocolor}{q} & \textcolor{gregoriocolor}{r} & \textcolor{gregoriocolor}{s} & \textcolor{gregoriocolor}{t} & \textcolor{gregoriocolor}{u} \\
 8 & 9 & 10 & 11 & 12 & 13 & 14 & 15 & 16 & 17 & 18 & 19 & 20 & 21 & 22 & 23 & 24 & 25 & 26 \\
\end{tabularx}
\vspace{0.5\baselineskip}
\noindent\begin{tabularx}{\linewidth}{*{12}{>{\centering\arraybackslash}X}}
 \textcolor{gregoriocolor}{A} & \textcolor{gregoriocolor}{B} & \textcolor{gregoriocolor}{C} & \textcolor{gregoriocolor}{D} & \textcolor{gregoriocolor}{E} & F & \textcolor{gregoriocolor}{F} & \textcolor{gregoriocolor}{G} & \textcolor{gregoriocolor}{H} & \textcolor{gregoriocolor}{M} & \textcolor{gregoriocolor}{N} & \textcolor{gregoriocolor}{P} \\
 27 & 28 & 29 & 1 & 2 & 3 & 2 & 3 & 4 & 5 & 6 & 7 \\
\end{tabularx}

\begin{paracol}{2}
\selectlanguage{latin}
\lettrine[lines=2]{M}{edioláni} pássio sancti Petri, ex Ordine Prædicatórum, 
 Mártyris, qui ab hæréticis, ob fidem cathólicam, interémptus est. 
 Ipsíus tamen festívitas recólitur tértio Kaléndas Maji.
\switchcolumn
\selectlanguage{english}
\lettrine[lines=2]{A}{t} Milan, the passion of St. Peter, a martyr belonging to the Order of 
 Preachers, who was slain by the heretics for his Catholic faith. His 
 feast, however, is kept on the 29th of April.
\switchcolumn*
\selectlanguage{latin}
Velehrádii, in Morávia, natális sancti Methódii, Epíscopi et Confessóris, 
 qui, una cum sancto Cyríllo, item Epíscopo et fratre suo, cujus natális 
 sextodécimo Kaléndas Mártii recensétur, multas Slávicas gentes earúmque 
 Reges ad fidem Christi perdúxit. Horum autem Sanctórum festum Nonis 
 Júlii celebrátur.
\switchcolumn
\selectlanguage{english}
In Moravia, the birthday of St. Methodius, bishop and confessor. 
 Together with his brother, the bishop St. Cyril, whose birthday was the 14th 
 of February, he converted many of the Slav races and their rulers to the 
 faith of Christ. Their feast is celebrated on the 7th day of July.
\switchcolumn*
\selectlanguage{latin}
In Macedónia sanctórum Mártyrum Timóthei et Diógenis.
\switchcolumn
\selectlanguage{english}
In Macedonia, the holy martyrs Timothy and Diogenes.
\switchcolumn*
\selectlanguage{latin}
In Pérside sanctórum centum vigínti Mártyrum.
\switchcolumn
\selectlanguage{english}
In Persia, one hundred and twenty holy martyrs.
\switchcolumn*
\selectlanguage{latin}
Ascalóne, in Palæstína, pássio sanctórum 
 Platónidis et aliórum duórum Mártyrum.
\switchcolumn
\selectlanguage{english}
At Ascalon in Palestine, the passion of St. Platonides and two other 
 martyrs.
\switchcolumn*
\selectlanguage{latin}
Carthágine sancti Marcellíni Mártyris, qui, ob cathólicæ 
 fidei defensiónem, ab hæréticis occísus est.
\switchcolumn
\selectlanguage{english}
At Carthage, St. Marcellin, who was slain by the heretics for defending the 
 Catholic faith.
\switchcolumn*
\selectlanguage{latin}
In Dánia sancti Guliélmi Abbátis, vita et miráculis clari.
\switchcolumn
\selectlanguage{english}
In Denmark, St. William, an abbot renowned for his saintly life and 
 miracles.
\switchcolumn*
\selectlanguage{latin}
\end{paracol}


% ---- martyrology/mart04/mart0407.htm
\needspace{10\baselineskip}
\begin{paracol}{2}
\selectlanguage{latin}
\begin{center}{\color{gregoriocolor} Séptimo Idus Aprílis. 
 Luna\dots\ }\end{center}
\switchcolumn
\selectlanguage{english}
\begin{center}{\color{gregoriocolor} The 7\textsuperscript{th} Day of April. The\dots\ Day of the Moon.}\end{center}
\end{paracol}

\noindent\begin{tabularx}{\linewidth}{*{19}{>{\centering\arraybackslash}X}}
 \textcolor{gregoriocolor}{a} & \textcolor{gregoriocolor}{b} & \textcolor{gregoriocolor}{c} & \textcolor{gregoriocolor}{d} & \textcolor{gregoriocolor}{e} & \textcolor{gregoriocolor}{f} & \textcolor{gregoriocolor}{g} & \textcolor{gregoriocolor}{h} & \textcolor{gregoriocolor}{i} & \textcolor{gregoriocolor}{k} & \textcolor{gregoriocolor}{l} & \textcolor{gregoriocolor}{m} & \textcolor{gregoriocolor}{n} & \textcolor{gregoriocolor}{p} & \textcolor{gregoriocolor}{q} & \textcolor{gregoriocolor}{r} & \textcolor{gregoriocolor}{s} & \textcolor{gregoriocolor}{t} & \textcolor{gregoriocolor}{u} \\
 9 & 10 & 11 & 12 & 13 & 14 & 15 & 16 & 17 & 18 & 19 & 20 & 21 & 22 & 23 & 24 & 25 & 26 & 27 \\
\end{tabularx}
\vspace{0.5\baselineskip}
\noindent\begin{tabularx}{\linewidth}{*{12}{>{\centering\arraybackslash}X}}
 \textcolor{gregoriocolor}{A} & \textcolor{gregoriocolor}{B} & \textcolor{gregoriocolor}{C} & \textcolor{gregoriocolor}{D} & \textcolor{gregoriocolor}{E} & F & \textcolor{gregoriocolor}{F} & \textcolor{gregoriocolor}{G} & \textcolor{gregoriocolor}{H} & \textcolor{gregoriocolor}{M} & \textcolor{gregoriocolor}{N} & \textcolor{gregoriocolor}{P} \\
 28 & 29 & 1 & 2 & 3 & 4 & 3 & 4 & 5 & 6 & 7 & 8 \\
\end{tabularx}

\begin{paracol}{2}
\selectlanguage{latin}
\lettrine[lines=2]{R}{otómagi} natális 
 sancti Joánnis Baptístæ de La Salle, Presbyteri et Confessóris, qui, in 
 erudiénda adolescéntia præsértim páupere excéllens, et de religióne 
 civilíque societáte præcláre méritus, Fratrum Scholárum Christianárum 
 Sodalitátem instítuit. Eum Pius Duodécimus, Póntifex Máximus, ómnium 
 Magistrórum púeris adolescentibúsque instituéndis præcípuum apud Deum 
 cæléstem Patrónum constítuit. Ipsíus tamen festum Idibus Maji 
 celebrátur.
\switchcolumn
\selectlanguage{english}
\lettrine[lines=2]{A}{t} Rouen, the birthday of St. John 
 Baptist de la Salle, priest and confessor. He was prominent in the 
 education of youth, especially those who were poor, for which he was 
 acclaimed both by religious and civil society. He was the founder of 
 the Society of the Brothers of the Christian Schools. Pius XII, 
 Supreme Pontiff, declared him patron of all those who teach children and 
 young people. His feast is celebrated on the 15th of May.
\switchcolumn*
\selectlanguage{latin}
In Africa item natális sanctórum Mártyrum Epiphánii Epíscopi, 
 Donáti, Rufíni et aliórum trédecim.
\switchcolumn
\selectlanguage{english}
In Africa, 
 the birthday of the holy martyrs Epiphanius bishop, Donatus, Rufinus and 
 thirteen others.
\switchcolumn*
\selectlanguage{latin}
Alexandríæ sancti Peléusii, Presbyteri, et Mártyris.
\switchcolumn
\selectlanguage{english}
At Alexandria, St. Peleusius, 
 priest and martyr.
\switchcolumn*
\selectlanguage{latin}
Synópe, in Ponto, sanctórum ducentórum Mártyrum.
\switchcolumn
\selectlanguage{english}
At Sinope, in Pontus, two hundred holy martyrs.
\switchcolumn*
\selectlanguage{latin}
In Cilícia sancti Calliópii Mártyris, qui, sub Maximiáno Præfécto, post 
 ália torménta, cápite in terram verso cruci affíxus, nóbili coróna martyrii 
 decorátus est.
\switchcolumn
\selectlanguage{english}
In Cilicia, under the prefect Maximus, St. Calliopius, martyr. 
 After undergoing other torments, he was fastened to a cross with his head 
 downward, and thus gained the noble crown of martyrdom.
\switchcolumn*
\selectlanguage{latin}
Nicomedíæ sancti Cyríaci et aliórum decem Mártyrum.
\switchcolumn
\selectlanguage{english}
At Nicomedia, St. Cyriacus and ten other martyrs.
\switchcolumn*
\selectlanguage{latin}
Verónæ sancti Saturníni, Epíscopi et Confessóris.
\switchcolumn
\selectlanguage{english}
At Verona, St. Saturninus, bishop and confessor.
\switchcolumn*
\selectlanguage{latin}
Romæ sancti Hegesíppi, qui vicínus Apostolórum tempóribus, Romam venit ad 
 Anicétum Pontíficem, ibíque mansit usque ad Eleuthérium, et Ecclesiasticórum 
 Actuum a passióne Dómini usque ad suam ætátem sermóne símplici téxuit 
 históriam; ut, quorum vitam sectabátur, dicéndi quoque exprímeret charactére.
\switchcolumn
\selectlanguage{english}
At Rome, St. Hegesippus, who lived close to the time of the apostles. 
 He came to Rome while Anicetus was pope, and remained until the time of 
 Eleutherius. He wrote a history of the Church, from the Passion of our 
 Lord to his own time, in a simple style, to make clear the character of 
 those whose life he imitated.
\switchcolumn*
\selectlanguage{latin}
In Syria sancti Aphraátis Anachorétæ, qui, Valéntis témpore, cathólicam 
 fidem virtúte miraculórum advérsus Ariános deféndit.
\switchcolumn
\selectlanguage{english}
In Syria, in the time of Valens, St. Aphraates, an anchoret, who defended 
 the Catholic faith against the Arians by the power of miracles.
\switchcolumn*
\selectlanguage{latin}
\end{paracol}


% ---- martyrology/mart04/mart0408.htm
\needspace{10\baselineskip}
\begin{paracol}{2}
\selectlanguage{latin}
\begin{center}{\color{gregoriocolor} Sexto Idus Aprílis. 
 Luna\dots\ }\end{center}
\switchcolumn
\selectlanguage{english}
\begin{center}{\color{gregoriocolor} The 8\textsuperscript{th} Day of April. The\dots\ Day of the Moon.}\end{center}
\end{paracol}

\noindent\begin{tabularx}{\linewidth}{*{19}{>{\centering\arraybackslash}X}}
 \textcolor{gregoriocolor}{a} & \textcolor{gregoriocolor}{b} & \textcolor{gregoriocolor}{c} & \textcolor{gregoriocolor}{d} & \textcolor{gregoriocolor}{e} & \textcolor{gregoriocolor}{f} & \textcolor{gregoriocolor}{g} & \textcolor{gregoriocolor}{h} & \textcolor{gregoriocolor}{i} & \textcolor{gregoriocolor}{k} & \textcolor{gregoriocolor}{l} & \textcolor{gregoriocolor}{m} & \textcolor{gregoriocolor}{n} & \textcolor{gregoriocolor}{p} & \textcolor{gregoriocolor}{q} & \textcolor{gregoriocolor}{r} & \textcolor{gregoriocolor}{s} & \textcolor{gregoriocolor}{t} & \textcolor{gregoriocolor}{u} \\
 10 & 11 & 12 & 13 & 14 & 15 & 16 & 17 & 18 & 19 & 20 & 21 & 22 & 23 & 24 & 25 & 26 & 27 & 28 \\
\end{tabularx}
\vspace{0.5\baselineskip}
\noindent\begin{tabularx}{\linewidth}{*{12}{>{\centering\arraybackslash}X}}
 \textcolor{gregoriocolor}{A} & \textcolor{gregoriocolor}{B} & \textcolor{gregoriocolor}{C} & \textcolor{gregoriocolor}{D} & \textcolor{gregoriocolor}{E} & F & \textcolor{gregoriocolor}{F} & \textcolor{gregoriocolor}{G} & \textcolor{gregoriocolor}{H} & \textcolor{gregoriocolor}{M} & \textcolor{gregoriocolor}{N} & \textcolor{gregoriocolor}{P} \\
 29 & 1 & 2 & 3 & 4 & 5 & 4 & 5 & 6 & 7 & 8 & 9 \\
\end{tabularx}

\begin{paracol}{2}
\selectlanguage{latin}
\lettrine[lines=2]{C}{ommemorátio} sanctórum 
 Herodiónis, Asyncriti et Phlegóntis, de quibus scribit beátus Paulus 
 Apóstolus in Epístola ad Romános.
\switchcolumn
\selectlanguage{english}
\lettrine[lines=2]{T}{he} commemoration of Saints 
 Herodian, Asyncritus, and Phlegon who are mentioned by blessed Paul the 
 Apostle in his Letter to the Romans.
\switchcolumn*
\selectlanguage{latin}
Alexandríæ sancti 
 Ædésii Mártyris, qui, sub Maximiáno Galério Imperatóre, cum esset beáti 
 Apphiáni frater, et ímpium Júdicem, quod Deo dicátas Vírgines lenónibus 
 tráderet, palam argúeret, idcírco, a milítibus tentus sævissimísque afféctus 
 supplíciis, in mare demérsus est pro Christo Dómino.
\switchcolumn
\selectlanguage{english}
At Alexandria, in the time of 
 Emperor Maximian Galerius, the martyr St. Aedesius, brother of the blessed 
 Apphian. Because he publicly reproved the wicked judge who delivered 
 to corruptors virgins consecrated to God, he was arrested by the soldiers, 
 exposed to the most severe torments, and thrown into the sea for the sake of 
 Christ our Lord.
\switchcolumn*
\selectlanguage{latin}
In Africa sanctórum 
 Mártyrum Januárii, Máximæ et Macáriæ.
\switchcolumn
\selectlanguage{english}
In Africa, the holy martyrs 
 Januarius, Maxima, and Macaria.
\switchcolumn*
\selectlanguage{latin}
Carthágine sanctæ 
 Concéssæ Mártyris.
\switchcolumn
\selectlanguage{english}
At Carthage, the martyr St. 
 Concessa.
\switchcolumn*
\selectlanguage{latin}
Apud Corínthum beáti 
 Dionysii Epíscopi, qui eruditióne et grátia, quam hábuit in verbo Dei, non 
 solum suæ civitátis et provínciæ pópulos, sed et aliárum provinciárum et 
 úrbium Epíscopos epístolis erudívit; Romanósque Pontífices 
 ádeo cóluit, ut 
 eórum epístolas públice légere in Ecclésia diébus Domínicis consuéverit. 
 Cláruit autem tempóribus Marci Antoníni Veri et Lúcii Aurélii Cómmodi.
\switchcolumn
\selectlanguage{english}
At Corinth, Bishop St. Denis, who 
 instructed not only the people of his own city and province by the learning 
 and charm with which he preached the word of God, but also the bishops of 
 other cities and provinces by the letters he wrote to them. His 
 devotion to the Roman Pontiffs was such that he was accustomed to read their 
 letters publicly in the church on Sundays. He lived in the time of 
 Marcus Antoninus Verus and Lucius Aurelius Commodus.
\switchcolumn*
\selectlanguage{latin}
Turónis, in Gállia, 
 sancti Perpétui Epíscopi, admirándæ sanctitátis viri.
\switchcolumn
\selectlanguage{english}
At Tours in France, the holy bishop 
 Perpetuus, a man of great sanctity.
\switchcolumn*
\selectlanguage{latin}
Ferentíni, in Hérnicis, 
 sancti Redémpti Epíscopi, cujus méminit beátus Gregórius Papa.
\switchcolumn
\selectlanguage{english}
At Ferentino in Campania, Bishop 
 St. Redemptus, who was mentioned by Pope St. Gregory.
\switchcolumn*
\selectlanguage{latin}
Apud Comum sancti 
 Amántii, Epíscopi et Confessóris.
\switchcolumn
\selectlanguage{english}
At Como, St. Amantius, bishop and 
 confessor.
\switchcolumn*
\selectlanguage{latin}
\end{paracol}


% ---- martyrology/mart04/mart0409.htm
\needspace{10\baselineskip}
\begin{paracol}{2}
\selectlanguage{latin}
\begin{center}{\color{gregoriocolor} Quinto Idus Aprílis. 
 Luna\dots\ }\end{center}
\switchcolumn
\selectlanguage{english}
\begin{center}{\color{gregoriocolor} The 
 9\textsuperscript{th} Day of April. The\dots\ Day of the Moon.}\end{center}
\end{paracol}

\noindent\begin{tabularx}{\linewidth}{*{19}{>{\centering\arraybackslash}X}}
 \textcolor{gregoriocolor}{a} & \textcolor{gregoriocolor}{b} & \textcolor{gregoriocolor}{c} & \textcolor{gregoriocolor}{d} & \textcolor{gregoriocolor}{e} & \textcolor{gregoriocolor}{f} & \textcolor{gregoriocolor}{g} & \textcolor{gregoriocolor}{h} & \textcolor{gregoriocolor}{i} & \textcolor{gregoriocolor}{k} & \textcolor{gregoriocolor}{l} & \textcolor{gregoriocolor}{m} & \textcolor{gregoriocolor}{n} & \textcolor{gregoriocolor}{p} & \textcolor{gregoriocolor}{q} & \textcolor{gregoriocolor}{r} & \textcolor{gregoriocolor}{s} & \textcolor{gregoriocolor}{t} & \textcolor{gregoriocolor}{u} \\
 11 & 12 & 13 & 14 & 15 & 16 & 17 & 18 & 19 & 20 & 21 & 22 & 23 & 24 & 25 & 26 & 27 & 28 & 29 \\
\end{tabularx}
\vspace{0.5\baselineskip}
\noindent\begin{tabularx}{\linewidth}{*{12}{>{\centering\arraybackslash}X}}
 \textcolor{gregoriocolor}{A} & \textcolor{gregoriocolor}{B} & \textcolor{gregoriocolor}{C} & \textcolor{gregoriocolor}{D} & \textcolor{gregoriocolor}{E} & F & \textcolor{gregoriocolor}{F} & \textcolor{gregoriocolor}{G} & \textcolor{gregoriocolor}{H} & \textcolor{gregoriocolor}{M} & \textcolor{gregoriocolor}{N} & \textcolor{gregoriocolor}{P} \\
 1 & 2 & 3 & 4 & 5 & 6 & 5 & 6 & 7 & 8 & 9 & 10 \\
\end{tabularx}

\begin{paracol}{2}
\selectlanguage{latin}
\lettrine[lines=2]{I}{n} Judǽa sanctæ Maríæ 
 Cléophæ, quam beátus Joánnes Evangelísta sorórem sanctíssimæ Dei Genitrícis 
 Maríæ núncupat, et cum hac simul juxta crucem Jesu stetísse narrat.
\switchcolumn
\selectlanguage{english}
\lettrine[lines=2]{I}{n} Judea, St. Mary Cleophas, whom 
 St. John the Evangelist calls the sister of the Blessed Virgin Mary, Mother 
 of God, and says that she stood at her side beneath the Cross of Jesus.
\switchcolumn*
\selectlanguage{latin}
Antiochíæ sancti 
 Próchori, qui fuit unus de septem primis Diáconis; et, fide ac miráculis 
 clarus, martyrio coronátus est.
\switchcolumn
\selectlanguage{english}
At Antioch, St. Prochorus who was 
 one of the first seven deacons. Renowned for faith and miracles, he 
 received the crown of martyrdom.
\switchcolumn*
\selectlanguage{latin}
Romæ natális sanctórum 
 Mártyrum Demétrii, Concéssi, Hilárii et Sociórum.
\switchcolumn
\selectlanguage{english}
At Rome, the birthday of the holy 
 martyrs Demetrius, Concessus, Hilary, and their companions.
\switchcolumn*
\selectlanguage{latin}
Cæsarǽæ, in Cappadócia, 
 sancti Eupsychii Mártyris, qui ob evérsum Fortúnæ fanum, sub Juliáno 
 Apóstata, martyrium consummávit.
\switchcolumn
\selectlanguage{english}
At Caesarea in Cappadocia, St. 
 Eupsychius, martyr, who was persecuted under Julian the Apostate for having 
 overthrown the temple of Fortune.
\switchcolumn*
\selectlanguage{latin}
In Africa sanctórum 
 Mártyrum Massylitanórum, in quorum die natáli sanctus Augustínus tractátum 
 hábuit.
\switchcolumn
\selectlanguage{english}
In Africa the holy Massylitan 
 Martyrs, on whose birthday was written a tract by St. Augustine.
\switchcolumn*
\selectlanguage{latin}
Sírmii pássio 
 sanctárum septem Vírginum et Mártyrum, quæ, dato simul prétio sánguinis, 
 vitam mercátæ sunt ætérnam.
\switchcolumn
\selectlanguage{english}
At Sirmio, seven holy virgins and 
 martyrs, who purchased eternal life together at the price of their own 
 blood.
\switchcolumn*
\selectlanguage{latin}
Amidæ, in Mesopotámia, sancti Acátii Epíscopi, qui pro rediméndis captívis étiam vasa Ecclésiæ 
 conflávit ac véndidit.
\switchcolumn
\selectlanguage{english}
At Amida in Mesopotamia, St. 
 Acatius, bishop, who even melted down and sold the sacred vessels in order 
 to ransom captives.
\switchcolumn*
\selectlanguage{latin}
Rotómagi sancti 
 Hugónis, Epíscopi et Confessóris.
\switchcolumn
\selectlanguage{english}
At Rouen, St. Hugh, bishop and 
 confessor.
\switchcolumn*
\selectlanguage{latin}
In civitáte Diénsi, in 
 Gállia, sancti Marcélli Epíscopi, miráculis clari.
\switchcolumn
\selectlanguage{english}
In the city of Die, in France, St. 
 Marcellus, bishop, celebrated for miracles.
\switchcolumn*
\selectlanguage{latin}
Móntibus, in Hannónia, 
 beátæ Waldetrúdis, vitæ sanctimónia et miráculis claræ.
\switchcolumn
\selectlanguage{english}
At Mons in Hainaut, blessed 
 Waltrude, renowned for holiness and miracles.
\switchcolumn*
\selectlanguage{latin}
Romæ Translátio 
 córporis sanctæ Mónicæ, matris beáti Augustíni Epíscopi; quod, ex Ostiis 
 Tiberínis, Martíno Quinto Summo Pontífice, in Urbem delátum, in Ecclésia 
 ejúsdem beáti Augustíni honorífice recónditum fuit.
\switchcolumn
\selectlanguage{english}
At Rome, the transferring of the 
 body of St. Monica, mother of the bishop St. Augustine. It was brought 
 from Ostia to Rome, under the Sovereign Pontiff, Martin V, and buried with 
 due honours in the church of St. Augustine.
\switchcolumn*
\selectlanguage{latin}
\end{paracol}


% ---- martyrology/mart04/mart0410.htm
\needspace{10\baselineskip}
\begin{paracol}{2}
\selectlanguage{latin}
\begin{center}{\color{gregoriocolor} Quarto Idus Aprílis. 
 Luna\dots\ }\end{center}
\switchcolumn
\selectlanguage{english}
\begin{center}{\color{gregoriocolor} The 10\textsuperscript{th} Day of April. The\dots\ Day of the Moon.}\end{center}
\end{paracol}

\noindent\begin{tabularx}{\linewidth}{*{19}{>{\centering\arraybackslash}X}}
 \textcolor{gregoriocolor}{a} & \textcolor{gregoriocolor}{b} & \textcolor{gregoriocolor}{c} & \textcolor{gregoriocolor}{d} & \textcolor{gregoriocolor}{e} & \textcolor{gregoriocolor}{f} & \textcolor{gregoriocolor}{g} & \textcolor{gregoriocolor}{h} & \textcolor{gregoriocolor}{i} & \textcolor{gregoriocolor}{k} & \textcolor{gregoriocolor}{l} & \textcolor{gregoriocolor}{m} & \textcolor{gregoriocolor}{n} & \textcolor{gregoriocolor}{p} & \textcolor{gregoriocolor}{q} & \textcolor{gregoriocolor}{r} & \textcolor{gregoriocolor}{s} & \textcolor{gregoriocolor}{t} & \textcolor{gregoriocolor}{u} \\
 12 & 13 & 14 & 15 & 16 & 17 & 18 & 19 & 20 & 21 & 22 & 23 & 24 & 25 & 26 & 27 & 28 & 29 & 1 \\
\end{tabularx}
\vspace{0.5\baselineskip}
\noindent\begin{tabularx}{\linewidth}{*{12}{>{\centering\arraybackslash}X}}
 \textcolor{gregoriocolor}{A} & \textcolor{gregoriocolor}{B} & \textcolor{gregoriocolor}{C} & \textcolor{gregoriocolor}{D} & \textcolor{gregoriocolor}{E} & F & \textcolor{gregoriocolor}{F} & \textcolor{gregoriocolor}{G} & \textcolor{gregoriocolor}{H} & \textcolor{gregoriocolor}{M} & \textcolor{gregoriocolor}{N} & \textcolor{gregoriocolor}{P} \\
 2 & 3 & 4 & 5 & 6 & 7 & 6 & 7 & 8 & 9 & 10 & 11 \\
\end{tabularx}

\begin{paracol}{2}
\selectlanguage{latin}
\lettrine[lines=2]{A}{pud} Babylónem sancti 
 Ezechiélis Prophétæ, qui, a Júdice pópuli Israël, quod eum de cultu idolórum 
 argúeret, interféctus, in sepúlcro Sem et Arpháxad, Abrahæ progenitórum, 
 sepúltus est; ad quod sepúlcrum, oratiónis causa, multi conflúere 
 consuevérunt.
\switchcolumn
\selectlanguage{english}
\lettrine[lines=2]{A}{t} Babylon, the prophet Ezechiel, 
 who was put to death by a judge of the people of Israel because he reproved 
 him for worshipping idols. He was buried in the sepulchre of Sem and 
 Arphaxad, ancestors of Abraham. Many people were in the habit of going 
 to his tomb to pray.
\switchcolumn*
\selectlanguage{latin}
Romæ natális 
 plurimórum sanctórum Mártyrum, quos sanctus Alexánder Papa, cum detinerétur 
 in cárcere, baptizávit. Hos autem omnes Aureliánus Præféctus, navi 
 vetústæ impósitos, in altum mare dedúci, et illic, ligátis ad colla 
 lapídibus, mergi præcépit.
\switchcolumn
\selectlanguage{english}
At Rome, the birthday of many holy 
 martyrs, whom Pope St. Alexander baptized while he was in prison. The 
 prefect Aurelian had them all put in an old ship, taken to the deep sea, and 
 drowned with stones tied to their necks.
\switchcolumn*
\selectlanguage{latin}
Alexandríæ sanctórum 
 Mártyrum Apollónii Presbyteri, et aliórum quinque, qui, in persecutióne 
 Maximiáni, in mare demérsi sunt.
\switchcolumn
\selectlanguage{english}
At Alexandria, during the 
 persecution of Maximian, the holy martyrs Apollonius, a priest, and five 
 others who were drowned in the sea.
\switchcolumn*
\selectlanguage{latin}
In Africa sanctórum 
 Mártyrum Teréntii, Africáni, Pompéji et Sociórum; qui, sub Décio Imperatóre 
 et Fortuniáno Præfécto, virgis cæsi, equúleis torti aliísque supplíciis 
 cruciáti, demum cápitis obtruncatióne martyrium complevérunt.
\switchcolumn
\selectlanguage{english}
In Africa, under Emperor Decius and 
 the prefect Fortunian, the holy martyrs Terence, Africanus, Pompey, and 
 their companions, who were scourged, racked and subjected to other torments. 
 Their martyrdom ended by beheading.
\switchcolumn*
\selectlanguage{latin}
Gandávi, in Flándria, 
 sancti Macárii, Epíscopi Antiochéni, virtútibus et miráculis clari.
\switchcolumn
\selectlanguage{english}
At Ghent in Flanders, St. Macarius, 
 bishop of Antioch, celebrated for virtues and miracles.
\switchcolumn*
\selectlanguage{latin}
Vallisoléti, in 
 Hispánia, sancti Michaélis de Sanctis, ex Ordine Discalceatórum sanctíssimæ 
 Trinitátis redemptiónis captivórum, Confessóris, innocéntia vitæ, admirábili 
 pæniténtia et caritáte in Deum exímii; quem Pius Nonus, Póntifex Máximus, 
 inter Sanctos rétulit.
\switchcolumn
\selectlanguage{english}
At Valladolid in Spain, St. Michael 
 of the Saints, confessor, of the Order of Discalced Trinitarians for the 
 Redemption of Captives, a man known for his upright life, his penitential 
 spirit, and his great love of God. He was placed on the roll of the 
 saints by Pope Pius IX.
\switchcolumn*
\selectlanguage{latin}
\end{paracol}


% ---- martyrology/mart04/mart0411.htm
\needspace{10\baselineskip}
\begin{paracol}{2}
\selectlanguage{latin}
\begin{center}{\color{gregoriocolor} Tértio Idus Aprílis. 
 Luna\dots\ }\end{center}
\switchcolumn
\selectlanguage{english}
\begin{center}{\color{gregoriocolor} The 11\textsuperscript{th} Day of April. The\dots\ Day of the Moon.}\end{center}
\end{paracol}

\noindent\begin{tabularx}{\linewidth}{*{19}{>{\centering\arraybackslash}X}}
 \textcolor{gregoriocolor}{a} & \textcolor{gregoriocolor}{b} & \textcolor{gregoriocolor}{c} & \textcolor{gregoriocolor}{d} & \textcolor{gregoriocolor}{e} & \textcolor{gregoriocolor}{f} & \textcolor{gregoriocolor}{g} & \textcolor{gregoriocolor}{h} & \textcolor{gregoriocolor}{i} & \textcolor{gregoriocolor}{k} & \textcolor{gregoriocolor}{l} & \textcolor{gregoriocolor}{m} & \textcolor{gregoriocolor}{n} & \textcolor{gregoriocolor}{p} & \textcolor{gregoriocolor}{q} & \textcolor{gregoriocolor}{r} & \textcolor{gregoriocolor}{s} & \textcolor{gregoriocolor}{t} & \textcolor{gregoriocolor}{u} \\
 13 & 14 & 15 & 16 & 17 & 18 & 19 & 20 & 21 & 22 & 23 & 24 & 25 & 26 & 27 & 28 & 29 & 1 & 2 \\
\end{tabularx}
\vspace{0.5\baselineskip}
\noindent\begin{tabularx}{\linewidth}{*{12}{>{\centering\arraybackslash}X}}
 \textcolor{gregoriocolor}{A} & \textcolor{gregoriocolor}{B} & \textcolor{gregoriocolor}{C} & \textcolor{gregoriocolor}{D} & \textcolor{gregoriocolor}{E} & F & \textcolor{gregoriocolor}{F} & \textcolor{gregoriocolor}{G} & \textcolor{gregoriocolor}{H} & \textcolor{gregoriocolor}{M} & \textcolor{gregoriocolor}{N} & \textcolor{gregoriocolor}{P} \\
 3 & 4 & 5 & 6 & 7 & 8 & 7 & 8 & 9 & 10 & 11 & 12 \\
\end{tabularx}

\begin{paracol}{2}
\selectlanguage{latin}
\lettrine[lines=2]{S}{ancti} Leónis Papæ 
 Primi, cognoménto Magni, Confessóris et Ecclésiæ Doctóris, cujus dies 
 natális recólitur quarto Idus Novémbris.
\switchcolumn
\selectlanguage{english}
\lettrine[lines=2]{S}{t.} Leo the First, pope and 
 confessor, who was surnamed the Great. His birthday falls on the 10th 
 of November.
\switchcolumn*
\selectlanguage{latin}
Pérgami, in Asia, 
 sancti Antípæ, testis fidélis, cujus méminit sanctus Joánnes in Apocalypsi. 
 Ipse autem Antípas, sub Domitiáno Imperatóre, in bovem æneum candéntem 
 conjéctus, martyrium consummávit.
\switchcolumn
\selectlanguage{english}
At Pergamum in Asia, the faithful 
 witness, St. Antipas, who was mentioned by St. John in the Apocalypse. 
 Under Emperor Domitian, he was enclosed in an ox made of brass that had been 
 heated to redness, and thus completed his martyrdom.
\switchcolumn*
\selectlanguage{latin}
Salónæ, in Dalmátia, 
 sanctórum Mártyrum Domniónis Epíscopi, cum milítibus octo.
\switchcolumn
\selectlanguage{english}
At Salona in Dalmatia, the holy 
 martyrs Domnio, bishop, and eight soldiers.
\switchcolumn*
\selectlanguage{latin}
Gortynæ, in Creta, 
 sancti Philíppi Epíscopi, vita et doctrína claríssimi, qui, tempóribus Marci 
 Antoníni Veri et Lúcii Aurélii Cómmodi, Ecclésiam sibi créditam regens, a 
 furóre Gentílium et ab hæreticórum insídiis eándem tutátus est.
\switchcolumn
\selectlanguage{english}
At Gortina in Crete, during the 
 reign of Marcus Antoninus Verus and Lucius Aurelius Commodus, St. Philip, 
 bishop, well known for his life and his teaching. He had defended the 
 Church entrusted to his care against the fury of the heathen and the snares 
 of the heretics.
\switchcolumn*
\selectlanguage{latin}
Nicomedíæ sancti 
 Eustórgii Presbyteri.
\switchcolumn
\selectlanguage{english}
At Nicomedia, the priest St. 
 Eustorgius.
\switchcolumn*
\selectlanguage{latin}
Spoléti sancti Isaac, 
 Mónachi et Confessóris, cujus virtútes sanctus Gregórius Papa commémorat.
\switchcolumn
\selectlanguage{english}
At Spoleto, St. Isaac, monk and 
 confessor, whose virtues are recorded by Pope St. Gregory.
\switchcolumn*
\selectlanguage{latin}
Apud Gazam in 
 Palæstína, sancti Barsanúphii Anachorétæ, sub Justiniáno Imperatóre.
\switchcolumn
\selectlanguage{english}
At Gaza in Palestine, in the time 
 of Emperor Justinian, St. Barsanuphius, an anchoret.
\switchcolumn*
\selectlanguage{latin}
Lucæ, in Etrúria, 
 sanctæ Gemmæ Galgáni, Vírginis, contemplatióne Domínicæ Passiónis et vitæ 
 sanctitátis mirábilis, quam Pius Papa Duodécimus in Sanctárum númerum 
 rétulit.
\switchcolumn
\selectlanguage{english}
At Luca in Etruria, St. Gemma 
 Galgani, virgin, renowned for her contemplation of the Passion of our Lord, 
 and for a life of holiness, and whom Pope Pius XII joined to the number of 
 the Saints.
\switchcolumn*
\selectlanguage{latin}
\end{paracol}


% ---- martyrology/mart04/mart0412.htm
\needspace{10\baselineskip}
\begin{paracol}{2}
\selectlanguage{latin}
\begin{center}{\color{gregoriocolor} Prídie Idus Aprílis. 
 Luna\dots\ }\end{center}
\switchcolumn
\selectlanguage{english}
\begin{center}{\color{gregoriocolor} The 12\textsuperscript{th} Day of April. The\dots\ Day of the Moon.}\end{center}
\end{paracol}

\noindent\begin{tabularx}{\linewidth}{*{19}{>{\centering\arraybackslash}X}}
 \textcolor{gregoriocolor}{a} & \textcolor{gregoriocolor}{b} & \textcolor{gregoriocolor}{c} & \textcolor{gregoriocolor}{d} & \textcolor{gregoriocolor}{e} & \textcolor{gregoriocolor}{f} & \textcolor{gregoriocolor}{g} & \textcolor{gregoriocolor}{h} & \textcolor{gregoriocolor}{i} & \textcolor{gregoriocolor}{k} & \textcolor{gregoriocolor}{l} & \textcolor{gregoriocolor}{m} & \textcolor{gregoriocolor}{n} & \textcolor{gregoriocolor}{p} & \textcolor{gregoriocolor}{q} & \textcolor{gregoriocolor}{r} & \textcolor{gregoriocolor}{s} & \textcolor{gregoriocolor}{t} & \textcolor{gregoriocolor}{u} \\
 14 & 15 & 16 & 17 & 18 & 19 & 20 & 21 & 22 & 23 & 24 & 25 & 26 & 27 & 28 & 29 & 1 & 2 & 3 \\
\end{tabularx}
\vspace{0.5\baselineskip}
\noindent\begin{tabularx}{\linewidth}{*{12}{>{\centering\arraybackslash}X}}
 \textcolor{gregoriocolor}{A} & \textcolor{gregoriocolor}{B} & \textcolor{gregoriocolor}{C} & \textcolor{gregoriocolor}{D} & \textcolor{gregoriocolor}{E} & F & \textcolor{gregoriocolor}{F} & \textcolor{gregoriocolor}{G} & \textcolor{gregoriocolor}{H} & \textcolor{gregoriocolor}{M} & \textcolor{gregoriocolor}{N} & \textcolor{gregoriocolor}{P} \\
 4 & 5 & 6 & 7 & 8 & 9 & 8 & 9 & 10 & 11 & 12 & 13 \\
\end{tabularx}

\begin{paracol}{2}
\selectlanguage{latin}
\lettrine[lines=2]{V}{erónæ} pássio sancti 
 Zenónis Epíscopi, qui inter persecutiónis procéllas eam Ecclésiam mira 
 constántia gubernávit, et, Galliéni témpore, martyrio coronátus est.
\switchcolumn
\selectlanguage{english}
\lettrine[lines=2]{A}{t} Verona, the passion of Bishop 
 St. Zeno, who governed that Church with great fortitude amid the storms of 
 persecution, and was crowned with martyrdom in the time of Gallienus.
\switchcolumn*
\selectlanguage{latin}
In Cappadócia sancti 
 Sabæ Gothi, qui, sub Valénte Imperatóre, cum Rex Gothórum Athanarícus 
 Christiános insequerétur, in flumen, post dira torménta, projéctus est; quo 
 étiam témpore (ut sanctus Augustínus scribit) quamplúrimi ex Gothis 
 orthodóxis coróna martyrii sunt insigníti.
\switchcolumn
\selectlanguage{english}
In Cappadocia, in the reign of 
 Emperor Valens, during the persecution raised against the Christians by 
 Atanaric, king of the Goths, St. Sabas, himself a Goth, who was cast into a 
 river after undergoing cruel torments. Many orthodox Goths, as St. 
 Augustine relates, received at that time the crown of martyrdom.
\switchcolumn*
\selectlanguage{latin}
Brácari, in Lusitánia, 
 sancti Victóris Mártyris, qui, adhuc catechúmenus, cum noluísset idólum 
 adoráre, et Christum Jesum magna constántia conféssus fuísset, ídeo, post 
 multa torménta, cápite abscísso, méruit próprio sánguine baptizári.
\switchcolumn
\selectlanguage{english}
At Braga in Portugal, the martyr 
 St. Victor. Although only a catechumen, he refused to adore an idol, 
 and confessed Jesus Christ with great constancy. After suffering many 
 tortures, he was beheaded, and thus merited to be baptized in his own blood.
\switchcolumn*
\selectlanguage{latin}
Firmi, in Picéno, 
 sanctæ Víssiæ, Vírginis et Mártyris.
\switchcolumn
\selectlanguage{english}
At Fermo, in Piceno, St. Vissia, 
 virgin and martyr.
\switchcolumn*
\selectlanguage{latin}
Romæ, via Aurélia, natális sancti Júlii Papæ Primi, qui advérsus Ariános pro fide cathólica 
 plúrimum laborávit, ac, multis præcláre gestis, sanctitáte célebris quiévit 
 in pace.
\switchcolumn
\selectlanguage{english}
At Rome, on the Aurelian Way, the 
 birthday of Pope St. Julius, who vigorously defended the Catholic faith 
 against the Arians. After a life of brilliant accomplishments, he 
 rested in peace, famed for his sanctity.
\switchcolumn*
\selectlanguage{latin}
Apud Vapíngum óppidum, 
 in Gállia, sancti Constantíni, Epíscopi et Confessóris.
\switchcolumn
\selectlanguage{english}
At the town of Gap in France, St. 
 Constantine, bishop and confessor.
\switchcolumn*
\selectlanguage{latin}
Papíæ sancti Damiáni 
 Epíscopi.
\switchcolumn
\selectlanguage{english}
At Pavia, Bishop St. Damian.
\switchcolumn*
\selectlanguage{latin}
\end{paracol}


% ---- martyrology/mart04/mart0413.htm
\needspace{10\baselineskip}
\begin{paracol}{2}
\selectlanguage{latin}
\begin{center}{\color{gregoriocolor} Idibus Aprílis. 
 Luna\dots\ }\end{center}
\switchcolumn
\selectlanguage{english}
\begin{center}{\color{gregoriocolor} The 13\textsuperscript{th} Day of April. The\dots\ Day of the Moon.}\end{center}
\end{paracol}

\noindent\begin{tabularx}{\linewidth}{*{19}{>{\centering\arraybackslash}X}}
 \textcolor{gregoriocolor}{a} & \textcolor{gregoriocolor}{b} & \textcolor{gregoriocolor}{c} & \textcolor{gregoriocolor}{d} & \textcolor{gregoriocolor}{e} & \textcolor{gregoriocolor}{f} & \textcolor{gregoriocolor}{g} & \textcolor{gregoriocolor}{h} & \textcolor{gregoriocolor}{i} & \textcolor{gregoriocolor}{k} & \textcolor{gregoriocolor}{l} & \textcolor{gregoriocolor}{m} & \textcolor{gregoriocolor}{n} & \textcolor{gregoriocolor}{p} & \textcolor{gregoriocolor}{q} & \textcolor{gregoriocolor}{r} & \textcolor{gregoriocolor}{s} & \textcolor{gregoriocolor}{t} & \textcolor{gregoriocolor}{u} \\
 15 & 16 & 17 & 18 & 19 & 20 & 21 & 22 & 23 & 24 & 25 & 26 & 27 & 28 & 29 & 1 & 2 & 3 & 4 \\
\end{tabularx}
\vspace{0.5\baselineskip}
\noindent\begin{tabularx}{\linewidth}{*{12}{>{\centering\arraybackslash}X}}
 \textcolor{gregoriocolor}{A} & \textcolor{gregoriocolor}{B} & \textcolor{gregoriocolor}{C} & \textcolor{gregoriocolor}{D} & \textcolor{gregoriocolor}{E} & F & \textcolor{gregoriocolor}{F} & \textcolor{gregoriocolor}{G} & \textcolor{gregoriocolor}{H} & \textcolor{gregoriocolor}{M} & \textcolor{gregoriocolor}{N} & \textcolor{gregoriocolor}{P} \\
 5 & 6 & 7 & 8 & 9 & 10 & 9 & 10 & 11 & 12 & 13 & 14 \\
\end{tabularx}

\begin{paracol}{2}
\selectlanguage{latin}
\lettrine[lines=2]{H}{íspali,} in Hispánia, sancti Hermenegíldi Mártyris, qui fuit fílius Leovigíldi, Regis Visigothórum 
 Ariáni; atque ob cathólicæ fidei confessiónem conjéctus in cárcerem, et, cum 
 in solemnitáte Pascháli Communiónem ab Epíscopo Ariáno accípere noluísset, 
 pérfidi patris jussu secúri percússus est, ac regnum cæléste pro terréno Rex 
 et Martyr intrávit.
\switchcolumn
\selectlanguage{english}
\lettrine[lines=2]{A}{t} Seville in Spain, St. 
 Hermenegild, son of Leovigild, Arian king of the Visigoths, who was 
 imprisoned for the confession of the Catholic faith. By order of his 
 wicked father he was beheaded because he had refused to receive communion 
 from an Arian bishop on the feast of Easter. Thus exchanging an 
 earthly for a heavenly kingdom, he entered the abode of the saints, both as 
 a king and as a martyr.
\switchcolumn*
\selectlanguage{latin}
Romæ, in persecutióne 
 Marci Antoníni Veri et Lúcii Aurélii Cómmodi, pássio sancti Justíni, 
 Philósophi et Mártyris; qui, cum secúndum librum pro religiónis nostræ 
 defensióne præfátis Imperatóribus porrexísset, eámque ibídem disputándo 
 strénue propugnásset, Crescéntis Cynici, cujus vitam et mores nefários 
 redargúerat, insídiis accusátus quod Christiánus esset, remuneratiónem 
 linguæ fidélis, martyrii munus accépit. Ipsíus tamen festum sequénti 
 die recólitur.
\switchcolumn
\selectlanguage{english}
At Rome, in the persecution of 
 Marcus Antoninus Verus and Lucius Aurelius Commodus, St. Justin, philosopher 
 and martyr. He had addressed to the emperors his second Apology in 
 defence of our religion, and upheld it by strong arguments. By the 
 intrigue of Crescens the Cynic, whose conduct and immorality he had 
 reproved, he was accused of professing Christianity, and thus he obtained 
 the reward of martyrdom in payment for his faithful confession. His 
 feast is kept on the following day.
\switchcolumn*
\selectlanguage{latin}
Pérgami, in Asia, in 
 eádem persecutióne, natális sanctórum Mártyrum Carpi, Thyatirénsis Epíscopi, 
 Pápyli Diáconi et Agathonícæ, ejúsdem Pápyli soróris et 
 óptimæ féminæ, atque 
 Agathadóri, eórum fámuli, aliorúmque multórum. Hi omnes, post vários 
 cruciátus, pro beátis confessiónibus martyrio coronáti sunt.
\switchcolumn
\selectlanguage{english}
At Pergamum in Asia, during the 
 same persecution, the birthday of the holy martyrs Carpus, bishop of 
 Thyatira, the deacon Papylus, and his sister Agathonica, an excellent woman, 
 Agathadorus, their servant, and many others. After many torments they 
 received their crowns of martyrdom for their worthy confessions.
\switchcolumn*
\selectlanguage{latin}
Doróstori, in Mysia 
 inferióre, pássio sanctórum Máximi, Quinctiliáni et Dadæ, in persecutióne 
 Diocletiáni.
\switchcolumn
\selectlanguage{english}
At Silistria in Bulgaria, the 
 passion of Saints Maximus, Quinctilian, and Dadas, during the persecution of 
 Diocletian.
\switchcolumn*
\selectlanguage{latin}
Ravénnæ sancti Ursi, 
 Epíscopi et Confessóris.
\switchcolumn
\selectlanguage{english}
At Ravenna, St. Ursus, bishop and 
 confessor.
\switchcolumn*
\selectlanguage{latin}
\end{paracol}


% ---- martyrology/mart04/mart0414.htm
\needspace{10\baselineskip}
\begin{paracol}{2}
\selectlanguage{latin}
\begin{center}{\color{gregoriocolor} Décimo octávo Kaléndas Maji. 
 Luna\dots\ }\end{center}
\switchcolumn
\selectlanguage{english}
\begin{center}{\color{gregoriocolor} The 14\textsuperscript{th} Day of April. The\dots\ Day of the Moon.}\end{center}
\end{paracol}

\noindent\begin{tabularx}{\linewidth}{*{19}{>{\centering\arraybackslash}X}}
 \textcolor{gregoriocolor}{a} & \textcolor{gregoriocolor}{b} & \textcolor{gregoriocolor}{c} & \textcolor{gregoriocolor}{d} & \textcolor{gregoriocolor}{e} & \textcolor{gregoriocolor}{f} & \textcolor{gregoriocolor}{g} & \textcolor{gregoriocolor}{h} & \textcolor{gregoriocolor}{i} & \textcolor{gregoriocolor}{k} & \textcolor{gregoriocolor}{l} & \textcolor{gregoriocolor}{m} & \textcolor{gregoriocolor}{n} & \textcolor{gregoriocolor}{p} & \textcolor{gregoriocolor}{q} & \textcolor{gregoriocolor}{r} & \textcolor{gregoriocolor}{s} & \textcolor{gregoriocolor}{t} & \textcolor{gregoriocolor}{u} \\
 16 & 17 & 18 & 19 & 20 & 21 & 22 & 23 & 24 & 25 & 26 & 27 & 28 & 29 & 1 & 2 & 3 & 4 & 5 \\
\end{tabularx}
\vspace{0.5\baselineskip}
\noindent\begin{tabularx}{\linewidth}{*{12}{>{\centering\arraybackslash}X}}
 \textcolor{gregoriocolor}{A} & \textcolor{gregoriocolor}{B} & \textcolor{gregoriocolor}{C} & \textcolor{gregoriocolor}{D} & \textcolor{gregoriocolor}{E} & F & \textcolor{gregoriocolor}{F} & \textcolor{gregoriocolor}{G} & \textcolor{gregoriocolor}{H} & \textcolor{gregoriocolor}{M} & \textcolor{gregoriocolor}{N} & \textcolor{gregoriocolor}{P} \\
 6 & 7 & 8 & 9 & 10 & 11 & 10 & 11 & 12 & 13 & 14 & 15 \\
\end{tabularx}

\begin{paracol}{2}
\selectlanguage{latin}
\lettrine[lines=2]{S}{ancti} Justíni, Philósophi et Mártyris, cujus memória 
 prídie hujus diéi recensétur.
\switchcolumn
\selectlanguage{english}
\lettrine[lines=2]{T}{he} feast of St. Justin, philosopher and martyr, who was 
 yesterday mentioned.
\switchcolumn*
\selectlanguage{latin}
Romæ, via Appia, 
 natális sanctórum Mártyrum Tibúrtii, Valeriáni et Máximi, sub Alexándro 
 Imperatóre et Almáchio Præfécto. Horum duo primi, beátæ Cæcíliæ 
 exhortatióne ad Christum convérsi et a sancto Urbáno Papa baptizáti, 
 póstmodum, ob fidei confessiónem, fústibus cæsi, gládio percússi sunt; 
 Máximus vero, Præfécti 
 cubiculárius, cum et ipse, eórum permótus constántia 
 et angélica visióne firmátus, in Christum credidísset, támdiu plumbátis 
 verberátus est, donec spíritum exhaláret.
\switchcolumn
\selectlanguage{english}
At Rome, on the Appian Way, the 
 birthday of the holy martyrs Tiburtius, Valerian, and Maximus, who suffered 
 in the time of Emperor Alexander and the prefect Almachius. The first 
 two were converted to Christ by the exhortations of blessed Cecilia, and 
 baptized by Pope St. Urban. They were beaten with clubs, then beheaded 
 for the sake of the true faith. Maximus, who had been the prefect's 
 chamberlain, was touched by their constancy, and confirmed by the vision of 
 an angel, believed in Christ, and was scourged with leaded whips until he 
 died.
\switchcolumn*
\selectlanguage{latin}
Interámnæ sancti 
 Próculi, Epíscopi et Mártyris.
\switchcolumn
\selectlanguage{english}
At Teramo, St. Proculus, bishop and 
 martyr.
\switchcolumn*
\selectlanguage{latin}
Eódem die sancti 
 Ardaliónis mimi, qui dum Sacris Christianórum in theatro illúderet, 
 derepénte mutátus est, et ea non solum verbis, sed étiam testimónio sui 
 sánguinis comprobávit.
\switchcolumn
\selectlanguage{english}
Also St. Ardalion, an actor. 
 One day in the theatre, while scoffing at the holy rites of the Christian 
 religion, he was suddenly converted and bore testimony to it, not only by 
 his words, but also with his blood.
\switchcolumn*
\selectlanguage{latin}
Interámnæ sanctæ 
 Domnínæ, Vírginis et Mártyris, cum Sóciis Virgínibus coronátæ.
\switchcolumn
\selectlanguage{english}
At Teramo, St. Domnina, virgin and 
 martyr, who received the crown with her virgin companions.
\switchcolumn*
\selectlanguage{latin}
Alexandríæ sanctæ 
 Thomáidis Mártyris, quæ a sócero, cujus impudícis nolúerat consentíre votis, 
 gládio percússa est atque in duas partes per médium discíssa.
\switchcolumn
\selectlanguage{english}
At Alexandria, St. Thomais, martyr. 
 Because she would not consent to the impure wishes of her father-in-law, she 
 was struck with a sword dividing her body from head to foot.
\switchcolumn*
\selectlanguage{latin}
Lugdúni, in Gállia, 
 sancti Lambérti, Epíscopi et Confessóris.
\switchcolumn
\selectlanguage{english}
At Lyons, in France, St. Lambert, bishop and 
 confessor.
\switchcolumn*
\selectlanguage{latin}
Alexandríæ sancti 
 Frontónis Abbátis, cujus vita sanctitáte et miráculis cláruit.
\switchcolumn
\selectlanguage{english}
At Alexandria, St. Fronto, an abbot 
 whose life was graced by sanctity and his miracles.
\switchcolumn*
\selectlanguage{latin}
Romæ sancti Abúndii, 
 Mansionárii Ecclésiæ sancti Petri.
\switchcolumn
\selectlanguage{english}
At Rome, St. Abundius, sacristan of 
 the church of St. Peter.
\switchcolumn*
\selectlanguage{latin}
\end{paracol}


% ---- martyrology/mart04/mart0415.htm
\needspace{10\baselineskip}
\begin{paracol}{2}
\selectlanguage{latin}
\begin{center}{\color{gregoriocolor} Décimo séptimo Kaléndas Maji. 
 Luna\dots\ }\end{center}
\switchcolumn
\selectlanguage{english}
\begin{center}{\color{gregoriocolor} The 15\textsuperscript{th} Day of April. The\dots\ Day of the Moon.}\end{center}
\end{paracol}

\noindent\begin{tabularx}{\linewidth}{*{19}{>{\centering\arraybackslash}X}}
 \textcolor{gregoriocolor}{a} & \textcolor{gregoriocolor}{b} & \textcolor{gregoriocolor}{c} & \textcolor{gregoriocolor}{d} & \textcolor{gregoriocolor}{e} & \textcolor{gregoriocolor}{f} & \textcolor{gregoriocolor}{g} & \textcolor{gregoriocolor}{h} & \textcolor{gregoriocolor}{i} & \textcolor{gregoriocolor}{k} & \textcolor{gregoriocolor}{l} & \textcolor{gregoriocolor}{m} & \textcolor{gregoriocolor}{n} & \textcolor{gregoriocolor}{p} & \textcolor{gregoriocolor}{q} & \textcolor{gregoriocolor}{r} & \textcolor{gregoriocolor}{s} & \textcolor{gregoriocolor}{t} & \textcolor{gregoriocolor}{u} \\
 17 & 18 & 19 & 20 & 21 & 22 & 23 & 24 & 25 & 26 & 27 & 28 & 29 & 1 & 2 & 3 & 4 & 5 & 6 \\
\end{tabularx}
\vspace{0.5\baselineskip}
\noindent\begin{tabularx}{\linewidth}{*{12}{>{\centering\arraybackslash}X}}
 \textcolor{gregoriocolor}{A} & \textcolor{gregoriocolor}{B} & \textcolor{gregoriocolor}{C} & \textcolor{gregoriocolor}{D} & \textcolor{gregoriocolor}{E} & F & \textcolor{gregoriocolor}{F} & \textcolor{gregoriocolor}{G} & \textcolor{gregoriocolor}{H} & \textcolor{gregoriocolor}{M} & \textcolor{gregoriocolor}{N} & \textcolor{gregoriocolor}{P} \\
 7 & 8 & 9 & 10 & 11 & 12 & 11 & 12 & 13 & 14 & 15 & 16 \\
\end{tabularx}

\begin{paracol}{2}
\selectlanguage{latin}
\lettrine[lines=2]{R}{omæ} sanctárum 
 Basilíssæ et Anastásiæ, nobílium feminárum, quæ, cum essent Apostolórum 
 discípulæ et constántes in fidei confessióne persísterent, sub Neróne 
 Imperatóre, lingua pedibúsque præcísis, percússæ gládio, martyrii corónam 
 adéptæ sunt.
\switchcolumn
\selectlanguage{english}
\lettrine[lines=2]{A}{t} Rome, the Saints Basilissa and 
 Anastasia, noble women who were disciples of the apostles. Because 
 they persevered courageously in the profession of their faith during 
 the time of the Emperor Nero, they had their tongues and feet cut off, were 
 put to the sword, and thus obtained the crown of martyrdom.
\switchcolumn*
\selectlanguage{latin}
Eódem die sanctórum 
 Mártyrum Marónis, Eutychétis et Victoríni; qui, primo cum beáta Flávia 
 Domitílla apud ínsulam Póntiam in Christi confessióne 
 éxsules, póstmodum, 
 sub Príncipe Nerva, liberáti, tandem, cum plúrimos ad fidem convertíssent, 
 in persecutióne Trajáni, a Valeriáno Júdice váriis pœnis jussi sunt 
 intérfici.
\switchcolumn
\selectlanguage{english}
The same day, the holy martyrs Maro, 
 Eutyches, and Victorinus, who, along with blessed Flavia Domitilla, had been 
 banished to the island of Pontia for the confession of Christ. Being 
 recalled in the reign of Nerva, and having converted many to the faith, they 
 were put to death in different ways by the judge Valerian, during the 
 persecution of Trajan.
\switchcolumn*
\selectlanguage{latin}
In Pérside sanctórum 
 Mártyrum Máximi et Olympíadis, qui, sub Décio Imperatóre, fústibus et 
 plumbátis cæsi sunt; et ad últimum cápita eórum sunt fústibus tunsa, donec 
 ambo emítterent spíritum.
\switchcolumn
\selectlanguage{english}
In Persia, in the reign of Emperor 
 Decius, the holy martyrs Maximus and Olympias, who were beaten with rods and 
 whips, and struck on their heads with clubs until they breathed no more.
\switchcolumn*
\selectlanguage{latin}
Ferentíni, in Hérnicis, 
 sancti Eutychii Mártyris.
\switchcolumn
\selectlanguage{english}
At Ferentino in Campania, the 
 martyr St. Eutychius.
\switchcolumn*
\selectlanguage{latin}
Myræ, in Lycia, sancti 
 Crescéntis, qui per ignem martyrium consummávit.
\switchcolumn
\selectlanguage{english}
At Myra in Lycia, St. Crescens, who 
 was martyred by fire.
\switchcolumn*
\selectlanguage{latin}
In Thrácia sanctórum 
 Mártyrum Theodóri et Pausilíppi, qui sub Hadriáno Imperatóre passi sunt.
\switchcolumn
\selectlanguage{english}
In Thrace, the holy martyrs 
 Theodorus and Pausilippus, who suffered under Emperor Hadrian.
\switchcolumn*
\selectlanguage{latin}
\end{paracol}


% ---- martyrology/mart04/mart0416.htm
\needspace{10\baselineskip}
\begin{paracol}{2}
\selectlanguage{latin}
\begin{center}{\color{gregoriocolor} Sextodécimo Kaléndas Maji. 
 Luna\dots\ }\end{center}
\switchcolumn
\selectlanguage{english}
\begin{center}{\color{gregoriocolor} The 16\textsuperscript{th} Day of April. The\dots\ Day of the Moon.}\end{center}
\end{paracol}

\noindent\begin{tabularx}{\linewidth}{*{19}{>{\centering\arraybackslash}X}}
 \textcolor{gregoriocolor}{a} & \textcolor{gregoriocolor}{b} & \textcolor{gregoriocolor}{c} & \textcolor{gregoriocolor}{d} & \textcolor{gregoriocolor}{e} & \textcolor{gregoriocolor}{f} & \textcolor{gregoriocolor}{g} & \textcolor{gregoriocolor}{h} & \textcolor{gregoriocolor}{i} & \textcolor{gregoriocolor}{k} & \textcolor{gregoriocolor}{l} & \textcolor{gregoriocolor}{m} & \textcolor{gregoriocolor}{n} & \textcolor{gregoriocolor}{p} & \textcolor{gregoriocolor}{q} & \textcolor{gregoriocolor}{r} & \textcolor{gregoriocolor}{s} & \textcolor{gregoriocolor}{t} & \textcolor{gregoriocolor}{u} \\
 18 & 19 & 20 & 21 & 22 & 23 & 24 & 25 & 26 & 27 & 28 & 29 & 1 & 2 & 3 & 4 & 5 & 6 & 7 \\
\end{tabularx}
\vspace{0.5\baselineskip}
\noindent\begin{tabularx}{\linewidth}{*{12}{>{\centering\arraybackslash}X}}
 \textcolor{gregoriocolor}{A} & \textcolor{gregoriocolor}{B} & \textcolor{gregoriocolor}{C} & \textcolor{gregoriocolor}{D} & \textcolor{gregoriocolor}{E} & F & \textcolor{gregoriocolor}{F} & \textcolor{gregoriocolor}{G} & \textcolor{gregoriocolor}{H} & \textcolor{gregoriocolor}{M} & \textcolor{gregoriocolor}{N} & \textcolor{gregoriocolor}{P} \\
 8 & 9 & 10 & 11 & 12 & 13 & 12 & 13 & 14 & 15 & 16 & 17 \\
\end{tabularx}

\begin{paracol}{2}
\selectlanguage{latin}
\lettrine[lines=2]{C}{orínthi} natális 
 sanctórum Mártyrum Callísti et Charísii, cum áliis septem, qui omnes, post 
 ália toleráta supplícia, in mare demérsi sunt.
\switchcolumn
\selectlanguage{english}
\lettrine[lines=2]{A}{t} Corinth, the birthday of the 
 holy martyrs Callistus and Charistius, with seven others, who were all cast 
 into the sea.
\switchcolumn*
\selectlanguage{latin}
Cæsaraugústæ, in 
 Hispánia, item natális sanctórum decem et octo Mártyrum, scílicet Optáti, 
 Lupérci, Succéssi, Martiális, Urbáni, Júliæ, Quinctiliáni, Públii, Frontónis, 
 Felícis, Cæciliáni, Evéntii, Primitívi, Apodémii, et aliórum quátuor qui 
 Saturníni vocáti esse referúntur. Hi omnes sub Daciáno, Hispaniárum 
 Præside, simul pœnis affécti atque interémpti sunt; quorum illústre 
 martyrium Prudéntius vérsibus exornávit.
\switchcolumn
\selectlanguage{english}
At Saragossa, in Spain, the 
 birthday of eighteen holy martyrs, Optatus, Lupercus, Successus, Martial, 
 Urban, Julia, Quinctilian, Publius, Fronto, Felix, Cecilian, Eventius, 
 Primitivus, Apodemius, and four others who are said to have been Saturninus. 
 They were all tortured and slain together under Dacian, governor of Spain. 
 The glory of their martyrdom has been celebrated in verse by Prudentius.
\switchcolumn*
\selectlanguage{latin}
In eádem civitáte 
 sanctórum Caji et Creméntii, qui secúndo conféssi et in Christi fide 
 perseverántes, martyrii cálicem gustavérunt.
\switchcolumn
\selectlanguage{english}
In the same city, the Saints Caius 
 and Crementius, who twice confessed the faith of Christ, and persevering in 
 it, drank of the chalice of martyrdom.
\switchcolumn*
\selectlanguage{latin}
Ibídem sancti Lambérti 
 Mártyris.
\switchcolumn
\selectlanguage{english}
In the same place, the martyr St. 
 Lambert.
\switchcolumn*
\selectlanguage{latin}
Item Cæsaraugústæ 
 sanctæ Encrátidis, Vírginis et Mártyris, quæ, laniáto córpore, mamílla 
 abscíssa et jécore avúlso, adhuc supérstes in cárcere inclúsa et reténta est, 
 donec ulcerátum ipsíus corpus putrésceret.
\switchcolumn
\selectlanguage{english}
Also at Saragossa, St. Encratis, 
 virgin and martyr, whose body was lacerated, her breasts cut away, and her 
 bowels torn out. Still alive after these torments, she was confined in 
 prison until her body, covered with wounds, began to decompose.
\switchcolumn*
\selectlanguage{latin}
Paléntiæ sancti 
 Turíbii, Epíscopi Asturicénsis, qui, ope sancti Leónis Papæ, Priscilliáni 
 hæresim ex Hispánia pénitus profligávit, clarúsque miráculis quiévit in 
 pace.
\switchcolumn
\selectlanguage{english}
At Palentia, St. Turibius, bishop 
 of Astorga. With the aid of Pope St. Leo, he drove out of Spain 
 completely the Priscillian heresy. He went to rest in the Lord with a 
 great renown for miracles.
\switchcolumn*
\selectlanguage{latin}
Brácari, in Lusitánia, 
 sancti Fructuósi Epíscopi.
\switchcolumn
\selectlanguage{english}
At Braga in Portugal, the bishop 
 St. Fructuosus.
\switchcolumn*
\selectlanguage{latin}
Sescíaci, in 
 Constantiénsi Gálliæ território, tránsitus sancti Patérni, Epíscopi 
 Abrincénsis et Confessóris.
\switchcolumn
\selectlanguage{english}
At Scicy, in the district of 
 Coutances in France, the death of St. Paternus, bishop of Avranches and 
 confessor.
\switchcolumn*
\selectlanguage{latin}
Romæ natális sancti 
 Benedícti-Joséphi Labre Confessóris, qui contémptu sui et extrémæ voluntáriæ 
 paupertátis laude exstitit insígnis.
\switchcolumn
\selectlanguage{english}
At Rome, the birthday of St. 
 Benedict Joseph Labre, confessor, who was famed for his contempt of self and 
 his great voluntary poverty.
\switchcolumn*
\selectlanguage{latin}
Apud Valencénas, in 
 Gállia, sancti Drogónis Confessóris.
\switchcolumn
\selectlanguage{english}
In Belgium, near Valenciennes, St. 
 Drogo, confessor.
\switchcolumn*
\selectlanguage{latin}
Nivérnis, in Gállia, 
 sanctæ Maríæ-Bernárdæ Soubirous, Vírginis, e Congregatióne Sorórum a 
 Caritáte et Institutióne Christiána, Lapúrdi, adhuc juvénculæ, iterátis 
 apparitiónibus Immaculátæ Dei Genitrícis Maríæ recreátæ; quam Pius Papa 
 Undécimus, inter sanctas Vírgines adscrípsit.
\switchcolumn
\selectlanguage{english}
In the city of Nevers in France, 
 St. Mary Bernard Soubirous of the Congregation of the Sisters of Charity, 
 also called the Christian Institute. She was favoured with frequent 
 apparitions and conversations at Lourdes with Mary Immaculate, the Mother of 
 God. In 1933 her name was added to the roll of holy virgins by Pope 
 Pius XI.
\switchcolumn*
\selectlanguage{latin}
Senis, in Túscia, 
 beáti Jóachim, ex Ordine Servórum beátæ Maríæ Vírginis.
\switchcolumn
\selectlanguage{english}
At Siena in Tuscany, blessed 
 Joachim of the Order of Servites of the Blessed Virgin Mary.
\switchcolumn*
\selectlanguage{latin}
\end{paracol}


% ---- martyrology/mart04/mart0417.htm
\needspace{10\baselineskip}
\begin{paracol}{2}
\selectlanguage{latin}
\begin{center}{\color{gregoriocolor} Quintodécimo Kaléndas Maji. 
 Luna\dots\ }\end{center}
\switchcolumn
\selectlanguage{english}
\begin{center}{\color{gregoriocolor} The 17\textsuperscript{th} Day of April. The\dots\ Day of the Moon.}\end{center}
\end{paracol}

\noindent\begin{tabularx}{\linewidth}{*{19}{>{\centering\arraybackslash}X}}
 \textcolor{gregoriocolor}{a} & \textcolor{gregoriocolor}{b} & \textcolor{gregoriocolor}{c} & \textcolor{gregoriocolor}{d} & \textcolor{gregoriocolor}{e} & \textcolor{gregoriocolor}{f} & \textcolor{gregoriocolor}{g} & \textcolor{gregoriocolor}{h} & \textcolor{gregoriocolor}{i} & \textcolor{gregoriocolor}{k} & \textcolor{gregoriocolor}{l} & \textcolor{gregoriocolor}{m} & \textcolor{gregoriocolor}{n} & \textcolor{gregoriocolor}{p} & \textcolor{gregoriocolor}{q} & \textcolor{gregoriocolor}{r} & \textcolor{gregoriocolor}{s} & \textcolor{gregoriocolor}{t} & \textcolor{gregoriocolor}{u} \\
 19 & 20 & 21 & 22 & 23 & 24 & 25 & 26 & 27 & 28 & 29 & 1 & 2 & 3 & 4 & 5 & 6 & 7 & 8 \\
\end{tabularx}
\vspace{0.5\baselineskip}
\noindent\begin{tabularx}{\linewidth}{*{12}{>{\centering\arraybackslash}X}}
 \textcolor{gregoriocolor}{A} & \textcolor{gregoriocolor}{B} & \textcolor{gregoriocolor}{C} & \textcolor{gregoriocolor}{D} & \textcolor{gregoriocolor}{E} & F & \textcolor{gregoriocolor}{F} & \textcolor{gregoriocolor}{G} & \textcolor{gregoriocolor}{H} & \textcolor{gregoriocolor}{M} & \textcolor{gregoriocolor}{N} & \textcolor{gregoriocolor}{P} \\
 9 & 10 & 11 & 12 & 13 & 14 & 13 & 14 & 15 & 16 & 17 & 18 \\
\end{tabularx}

\begin{paracol}{2}
\selectlanguage{latin}
\lettrine[lines=2]{R}{omæ} sancti Anicéti, 
 Papæ et Mártyris; qui martyrii palmam in persecutióne Marci Aurélii Antoníni 
 et Lúcii Veri accépit.
\switchcolumn
\selectlanguage{english}
\lettrine[lines=2]{A}{t} Rome, St. Anicetus, pope and 
 martyr, who received the palm of martyrdom in the persecution of Marcus 
 Aurelius Antoninus and Lucius Verus.
\switchcolumn*
\selectlanguage{latin}
Córdubæ, in Hispánia, sanctórum Mártyrum Elíæ Presbyteri, Pauli et Isidóri Monachórum, qui, in 
 persecutióne Arábica, ob Christiánæ fidei professiónem, interémpti sunt.
\switchcolumn
\selectlanguage{english}
At Cordova in Spain, the holy 
 martyrs Elias, a priest, and the monks Paul and Isidore, who were slain in 
 the Arab persecution for the profession of the Christian faith.
\switchcolumn*
\selectlanguage{latin}
Antiochíæ sanctórum 
 Mártyrum Petri Diáconi, et Hermógenis, qui ejúsdem Petri erat miníster.
\switchcolumn
\selectlanguage{english}
At Antioch, the holy martyrs Peter, 
 a deacon, and Hermogenes, who was his servant.
\switchcolumn*
\selectlanguage{latin}
In Africa natális 
 beáti Mappálici Mártyris, qui (ut scribit sanctus Cypriánus in epístola ad 
 Mártyres et Confessóres), cum áliis plúribus, martyrio coronátus est.
\switchcolumn
\selectlanguage{english}
In Africa, the birthday of blessed 
 Mappalicus, martyr. St. Cyprian relates in his Epistle to the Martyrs 
 and Confessors that he, along with many others, was crowned with martyrdom.
\switchcolumn*
\selectlanguage{latin}
Ibídem sanctórum 
 Mártyrum Fortunáti et Marciáni.
\switchcolumn
\selectlanguage{english}
In the same place, the holy martyrs 
 Fortunatus and Marcian.
\switchcolumn*
\selectlanguage{latin}
Viénnæ, in Gállia, 
 sancti Pantágathi Epíscopi.
\switchcolumn
\selectlanguage{english}
At Vienne in France, Bishop St. 
 Pantagathus.
\switchcolumn*
\selectlanguage{latin}
Dertónæ sancti 
 Innocéntii, Epíscopi et Confessóris.
\switchcolumn
\selectlanguage{english}
At Tortona, St. Innocent, bishop 
 and confessor.
\switchcolumn*
\selectlanguage{latin}
Cistércii, in Gállia, 
 sancti Stéphani Abbátis, qui primus erémum Cisterciénsem incóluit, et 
 sanctum Bernárdum, cum sóciis ad se veniéntem, lætus excépit.
\switchcolumn
\selectlanguage{english}
At Citeaux in France, St. Stephen Harding, 
 abbot, who was first to live in the Cistercian desert and who joyfully 
 welcomed St. Bernard and his companions when they came to him.
\switchcolumn*
\selectlanguage{latin}
In monastério Casæ 
 Dei, Claromontánæ diœcésis, in Gállia, sancti Robérti Confessóris, qui 
 ejúsdem monastérii cónditor et primus Abbas éxstitit.
\switchcolumn
\selectlanguage{english}
In the monastery of Chaise-Dieu, in 
 the diocese of Clermont, St. Robert, confessor, the founder and first abbot 
 of the monastery.
\switchcolumn*
\selectlanguage{latin}
\end{paracol}


% ---- martyrology/mart04/mart0418.htm
\needspace{10\baselineskip}
\begin{paracol}{2}
\selectlanguage{latin}
\begin{center}{\color{gregoriocolor} Quartodécimo Kaléndas Maji. 
 Luna\dots\ }\end{center}
\switchcolumn
\selectlanguage{english}
\begin{center}{\color{gregoriocolor} The 18\textsuperscript{th} Day of April. The\dots\ Day of the Moon.}\end{center}
\end{paracol}

\noindent\begin{tabularx}{\linewidth}{*{19}{>{\centering\arraybackslash}X}}
 \textcolor{gregoriocolor}{a} & \textcolor{gregoriocolor}{b} & \textcolor{gregoriocolor}{c} & \textcolor{gregoriocolor}{d} & \textcolor{gregoriocolor}{e} & \textcolor{gregoriocolor}{f} & \textcolor{gregoriocolor}{g} & \textcolor{gregoriocolor}{h} & \textcolor{gregoriocolor}{i} & \textcolor{gregoriocolor}{k} & \textcolor{gregoriocolor}{l} & \textcolor{gregoriocolor}{m} & \textcolor{gregoriocolor}{n} & \textcolor{gregoriocolor}{p} & \textcolor{gregoriocolor}{q} & \textcolor{gregoriocolor}{r} & \textcolor{gregoriocolor}{s} & \textcolor{gregoriocolor}{t} & \textcolor{gregoriocolor}{u} \\
 20 & 21 & 22 & 23 & 24 & 25 & 26 & 27 & 28 & 29 & 1 & 2 & 3 & 4 & 5 & 6 & 7 & 8 & 9 \\
\end{tabularx}
\vspace{0.5\baselineskip}
\noindent\begin{tabularx}{\linewidth}{*{12}{>{\centering\arraybackslash}X}}
 \textcolor{gregoriocolor}{A} & \textcolor{gregoriocolor}{B} & \textcolor{gregoriocolor}{C} & \textcolor{gregoriocolor}{D} & \textcolor{gregoriocolor}{E} & F & \textcolor{gregoriocolor}{F} & \textcolor{gregoriocolor}{G} & \textcolor{gregoriocolor}{H} & \textcolor{gregoriocolor}{M} & \textcolor{gregoriocolor}{N} & \textcolor{gregoriocolor}{P} \\
 10 & 11 & 12 & 13 & 14 & 15 & 14 & 15 & 16 & 17 & 18 & 19 \\
\end{tabularx}

\begin{paracol}{2}
\selectlanguage{latin}
\lettrine[lines=2]{A}{pud} montem Senárium, 
 in Etrúria, natális sancti Amidǽi Confessóris, e septem Fundatóribus Ordinis 
 Servórum beátæ Maríæ Vírginis, flagrantíssima in Deum caritáte præclári. 
 Ipsíus tamen ac Sociórum festum prídie Idus Februárii celebrátur.
\switchcolumn
\selectlanguage{english}
\lettrine[lines=2]{O}{n} Mount Senario in Tuscany, St. 
 Amadeo, confessor, one of the seven founders of the Order of Servites of the 
 Blessed Virgin Mary, famous for his ardent love for God. His feast, 
 together with that of his companions, is kept on the 12th of February.
\switchcolumn*
\selectlanguage{latin}
Romæ beáti Apollónii 
 Senatóris, qui, sub Cómmodo Príncipe et Perénnio Præfécto, a servo próditus 
 quod Christiánus esset, jussúsque ut ratiónem fídei suæ rédderet, insígne volúmen compósuit, quod in Senátu legit; et nihilóminus pro Christo, 
 senténtia Senátus, cápite truncátus est.
\switchcolumn
\selectlanguage{english}
At Rome, blessed Apollonius, a 
 senator under Emperor Commodus and the prefect Perennius. He was 
 denounced as a Christian by one of his slaves, and being commanded to give 
 an account of his faith, he composed an able work which he read in the 
 Senate. He was nevertheless beheaded for Christ by their sentence.
\switchcolumn*
\selectlanguage{latin}
Messánæ, in Sicília, 
 natális sanctórum Mártyrum Eleuthérii, Epíscopi Illyrici, et Anthíæ matris. 
 Ipse, cum esset vitæ sanctimónia et miraculórum virtúte illústris, sub 
 Hadriáno Príncipe, lectum férreum ignítum, cratículam et sartáginem, óleo et 
 pice ac resína fervéntem, súperans, projéctus quoque leónibus sed ab illis 
 nil læsus, novíssime simul cum matre jugulátur.
\switchcolumn
\selectlanguage{english}
At Messina in Sicily, the birthday 
 of the holy martyrs Eleutherius, bishop of Illyria, and Anthia, his mother. 
 He was famous for holiness of life and the power of miracles. During 
 the reign of Hadrian, he was placed on a bed of red-hot iron, on a gridiron, 
 in a vessel filled with boiling oil, pitch, and resin, and also cast to the 
 lions; but remaining unhurt through all of this, they finally cut his throat 
 with a sword. His mother suffered the same torments.
\switchcolumn*
\selectlanguage{latin}
Córdubæ, in Hispánia, sancti Perfécti, Presbyteri et Mártyris; qui a Mauris, eo quod inveherétur 
 in Mahumétis sectam et fírmiter Christi fidem profiterétur, gládio 
 trucidátus est.
\switchcolumn
\selectlanguage{english}
At Cordova, St. Perfectus, priest 
 and martyr, who was slain with the sword by the Moors, because he argued 
 against the sect of Mohammed and firmly insisted on the Catholic faith.
\switchcolumn*
\selectlanguage{latin}
Messánæ, in Sicília, 
 sancti Corébi Præfécti, qui, per sanctum Eleuthérium convérsus ad fidem, 
 gládio percússus est.
\switchcolumn
\selectlanguage{english}
At Messina in Sicily, St. Corebus, 
 the prefect, who was converted to the faith by St. Eleutherius, and died by 
 the sword.
\switchcolumn*
\selectlanguage{latin}
Bríxiæ sancti Calóceri 
 Mártyris, qui, a sanctis Faustíno et Jovíta convérsus ad Christum, sub 
 Hadriáno Príncipe gloriósum confessiónis certámen complévit.
\switchcolumn
\selectlanguage{english}
At Brescia, the martyr St. 
 Calocerus, who was converted to Christ by Saints Faustinus and Jovita, and 
 who gloriously triumphed in the test of his confession, in the time of 
 Hadrian.
\switchcolumn*
\selectlanguage{latin}
Medioláni sancti 
 Galdíni, Cardinális et ejúsdem civitátis Epíscopi, qui, concióne advérsus 
 hæréticos expléta, spíritum Deo réddidit.
\switchcolumn
\selectlanguage{english}
At Milan, St. Galdino, cardinal 
 bishop of that city, who at the very end of a sermon against heretics, gave 
 up his soul to God.
\switchcolumn*
\selectlanguage{latin}
\end{paracol}


% ---- martyrology/mart04/mart0419.htm
\needspace{10\baselineskip}
\begin{paracol}{2}
\selectlanguage{latin}
\begin{center}{\color{gregoriocolor} Tertiodécimo Kaléndas Maji. 
 Luna\dots\ }\end{center}
\switchcolumn
\selectlanguage{english}
\begin{center}{\color{gregoriocolor} The 19\textsuperscript{th} Day of April. The\dots\ Day of the Moon.}\end{center}
\end{paracol}

\noindent\begin{tabularx}{\linewidth}{*{19}{>{\centering\arraybackslash}X}}
 \textcolor{gregoriocolor}{a} & \textcolor{gregoriocolor}{b} & \textcolor{gregoriocolor}{c} & \textcolor{gregoriocolor}{d} & \textcolor{gregoriocolor}{e} & \textcolor{gregoriocolor}{f} & \textcolor{gregoriocolor}{g} & \textcolor{gregoriocolor}{h} & \textcolor{gregoriocolor}{i} & \textcolor{gregoriocolor}{k} & \textcolor{gregoriocolor}{l} & \textcolor{gregoriocolor}{m} & \textcolor{gregoriocolor}{n} & \textcolor{gregoriocolor}{p} & \textcolor{gregoriocolor}{q} & \textcolor{gregoriocolor}{r} & \textcolor{gregoriocolor}{s} & \textcolor{gregoriocolor}{t} & \textcolor{gregoriocolor}{u} \\
 21 & 22 & 23 & 24 & 25 & 26 & 27 & 28 & 29 & 1 & 2 & 3 & 4 & 5 & 6 & 7 & 8 & 9 & 10 \\
\end{tabularx}
\vspace{0.5\baselineskip}
\noindent\begin{tabularx}{\linewidth}{*{12}{>{\centering\arraybackslash}X}}
 \textcolor{gregoriocolor}{A} & \textcolor{gregoriocolor}{B} & \textcolor{gregoriocolor}{C} & \textcolor{gregoriocolor}{D} & \textcolor{gregoriocolor}{E} & F & \textcolor{gregoriocolor}{F} & \textcolor{gregoriocolor}{G} & \textcolor{gregoriocolor}{H} & \textcolor{gregoriocolor}{M} & \textcolor{gregoriocolor}{N} & \textcolor{gregoriocolor}{P} \\
 11 & 12 & 13 & 14 & 15 & 16 & 15 & 16 & 17 & 18 & 19 & 20 \\
\end{tabularx}

\begin{paracol}{2}
\selectlanguage{latin}
\lettrine[lines=2]{C}{orínthi} natális 
 sancti Timónis, qui fuit unus de septem primis Diáconis. Hic primo 
 apud Berœam Doctor resédit, ac deínde, verbum Dómini disséminans, venit 
 Corínthum; ibíque, a Judǽis et Græcis (ut tráditur) injéctus flammis, sed 
 nihil læsus, demum, cruci affíxus, martyrium suum implévit.
\switchcolumn
\selectlanguage{english}
\lettrine[lines=2]{A}{t} Corinth, the birthday of St. 
 Timon, one of the first seven deacons, who was first a teacher at Berea. 
 Afterwards, while preaching the word of the Lord at Corinth, he was 
 delivered to the flames by the Jews and the Greeks, but remaining uninjured, 
 he ended his martyrdom by crucifixion.
\switchcolumn*
\selectlanguage{latin}
Cantuáriæ, in Anglia, 
 sancti Elphégi, Epíscopi et Mártyris.
\switchcolumn
\selectlanguage{english}
At Canterbury in England, St. 
 Alphege, bishop and martyr.
\switchcolumn*
\selectlanguage{latin}
Melitínæ, in Arménia, sanctórum Mártyrum Hermógenis, Caji, Expedíti, Aristónici, Rufi et Galátæ, 
 qui omnes una die sunt coronáti.
\switchcolumn
\selectlanguage{english}
At Melitine in Armenia, the holy 
 martyrs Hermogenes, Caius, Expeditus, Aristonicus, Rufus, and Galatas, all 
 crowned on the same day.
\switchcolumn*
\selectlanguage{latin}
Caucolíberi, in 
 Hispánia Tarraconénsi, pássio sancti Vincéntii Mártyris.
\switchcolumn
\selectlanguage{english}
At Collioure in Spain, the martyr 
 St. Vincent.
\switchcolumn*
\selectlanguage{latin}
Eódem die sanctórum 
 Mártyrum Socrátis et Dionysii, qui lánceis confóssi sunt.
\switchcolumn
\selectlanguage{english}
On the same day, the holy martyrs 
 Socrates and Denis, who were killed with spears.
\switchcolumn*
\selectlanguage{latin}
Hierosólymis sancti 
 Paphnútii Mártyris.
\switchcolumn
\selectlanguage{english}
At Jerusalem, the martyr St. 
 Paphnutius.
\switchcolumn*
\selectlanguage{latin}
Romæ sancti Leónis 
 Papæ Noni, virtútum et miraculórum laude insígnis.
\switchcolumn
\selectlanguage{english}
At Rome, Pope St. Leo IX, 
 illustrious for his virtues and his miracles.
\switchcolumn*
\selectlanguage{latin}
Antiochíæ Pisídiæ 
 sancti Geórgii Epíscopi, qui, ob sanctárum Imáginum cultum, exsul occúbuit.
\switchcolumn
\selectlanguage{english}
At Antioch in Pisidia, St. George, 
 a bishop, who died in exile for the veneration of sacred images.
\switchcolumn*
\selectlanguage{latin}
In cœnóbio Laubiénsi, 
 in Bélgio, sancti Ursmári Epíscopi.
\switchcolumn
\selectlanguage{english}
In the monastery of Lobbes in Belgium, the 
 bishop St. Ursmar.
\switchcolumn*
\selectlanguage{latin}
Floréntiæ sancti 
 Crescéntii Confessóris, qui fuit discípulus beáti Zenóbii Epíscopi.
\switchcolumn
\selectlanguage{english}
At Florence, St. Crescent, 
 confessor, a disciple of the blessed Bishop Zenobius.
\switchcolumn*
\selectlanguage{latin}
\end{paracol}


% ---- martyrology/mart04/mart0420.htm
\needspace{10\baselineskip}
\begin{paracol}{2}
\selectlanguage{latin}
\begin{center}{\color{gregoriocolor} Duodécimo Kaléndas Maji. 
 Luna\dots\ }\end{center}
\switchcolumn
\selectlanguage{english}
\begin{center}{\color{gregoriocolor} The 
 20\textsuperscript{th} Day of April. The\dots\ Day of the Moon.}\end{center}
\end{paracol}

\noindent\begin{tabularx}{\linewidth}{*{19}{>{\centering\arraybackslash}X}}
 \textcolor{gregoriocolor}{a} & \textcolor{gregoriocolor}{b} & \textcolor{gregoriocolor}{c} & \textcolor{gregoriocolor}{d} & \textcolor{gregoriocolor}{e} & \textcolor{gregoriocolor}{f} & \textcolor{gregoriocolor}{g} & \textcolor{gregoriocolor}{h} & \textcolor{gregoriocolor}{i} & \textcolor{gregoriocolor}{k} & \textcolor{gregoriocolor}{l} & \textcolor{gregoriocolor}{m} & \textcolor{gregoriocolor}{n} & \textcolor{gregoriocolor}{p} & \textcolor{gregoriocolor}{q} & \textcolor{gregoriocolor}{r} & \textcolor{gregoriocolor}{s} & \textcolor{gregoriocolor}{t} & \textcolor{gregoriocolor}{u} \\
 22 & 23 & 24 & 25 & 26 & 27 & 28 & 29 & 1 & 2 & 3 & 4 & 5 & 6 & 7 & 8 & 9 & 10 & 11 \\
\end{tabularx}
\vspace{0.5\baselineskip}
\noindent\begin{tabularx}{\linewidth}{*{12}{>{\centering\arraybackslash}X}}
 \textcolor{gregoriocolor}{A} & \textcolor{gregoriocolor}{B} & \textcolor{gregoriocolor}{C} & \textcolor{gregoriocolor}{D} & \textcolor{gregoriocolor}{E} & F & \textcolor{gregoriocolor}{F} & \textcolor{gregoriocolor}{G} & \textcolor{gregoriocolor}{H} & \textcolor{gregoriocolor}{M} & \textcolor{gregoriocolor}{N} & \textcolor{gregoriocolor}{P} \\
 12 & 13 & 14 & 15 & 16 & 17 & 16 & 17 & 18 & 19 & 20 & 21 \\
\end{tabularx}

\begin{paracol}{2}
\selectlanguage{latin}
\lettrine[lines=2]{R}{omæ} sanctórum 
 Mártyrum Sulpícii et Serviliáni, qui, prædicatióne et miráculis beátæ 
 Domitíllæ Vírginis ad Christi fidem convérsi, ambo, cum nollent idólis 
 immoláre, in persecutióne Trajáni, a Præfécto Urbis Aniáno sunt cápite cæsi.
\switchcolumn
\selectlanguage{english}
\lettrine[lines=2]{A}{t} Rome, the holy martyrs Sulpicius 
 and Servilian, who were converted to the faith of Christ by the speeches and 
 the miracles of the holy virgin Domitilla. Because they refused to 
 sacrifice to the idols, they were beheaded by Anian, prefect of the city, in 
 the persecution of Trajan.
\switchcolumn*
\selectlanguage{latin}
Nicomedíæ sanctórum 
 Mártyrum Victóris, Zótici, Zenónis, Acíndyni, Cæsárei, Severiáni, 
 Chrysóphori, Theónæ et Antoníni, qui, sub Diocletiáno Imperatóre, passióne 
 ac signis beáti Geórgii ad Christum convérsi sunt, et ob intrépidam fídei 
 confessiónem, várie tentáti, martyrium complevérunt.
\switchcolumn
\selectlanguage{english}
At Nicomedia, the holy martyrs 
 Victor, Zoticus, Zeno, Acindynus, Caesareus, Severian, Chrysophorus, Theonas, 
 and Antonine. They were converted to Christ by the miracles and the 
 martyrdom of St. George, and because of their own dauntless confession of 
 the faith, they were tortured in various ways under the Emperor Diocletian, 
 and thus completed their martyrdom.
\switchcolumn*
\selectlanguage{latin}
Tomis, in Scythia, 
 sancti Theótimi Epíscopi, quem, ob insígnem ipsíus sanctitátem atque 
 mirácula, étiam infidéles bárbari veneráti sunt.
\switchcolumn
\selectlanguage{english}
At Tomis in Scythia, Bishop St. 
 Theotimus, whose great sanctity and miracles procured him the respect even 
 of unbelieving barbarians.
\switchcolumn*
\selectlanguage{latin}
Ebredúni, in Gálliis, 
 sancti Marcellíni, qui fuit primus ejúsdem urbis Epíscopus. Hic, Dei 
 mónitu, cum sanctis Sóciis Vincéntio et Domníno, ex Africa venit, et máximam 
 Alpium maritimárum partem verbo et signis admirándis, quibus usque hódie 
 refúlget, ad Christi fidem convértit.
\switchcolumn
\selectlanguage{english}
At Embrun in France, St. Marcellin, 
 first bishop of that city. By divine inspiration he came from Africa 
 with his holy companions Vincent and Domninus, and converted the greater 
 portion of the inhabitants of the Maritime Alps by his preaching, and by the 
 wonderful prodigies which he still continues to work.
\switchcolumn*
\selectlanguage{latin}
Antisiodóri sancti 
 Marciáni Presbyteri.
\switchcolumn
\selectlanguage{english}
At Auxerre, the priest St. Marcian.
\switchcolumn*
\selectlanguage{latin}
Apud Constantinópolim 
 sancti Theodóri Confessóris, ab áspera cilícii veste, qua tegebátur, 
 cognoménto Tríchinas, qui multis virtútibus, præsértim advérsus dæmones, 
 cláruit; ex cujus córpore scatúriens unguéntum ægrótis sanitátem impértit.
\switchcolumn
\selectlanguage{english}
At Constantinople, St. Theodore, 
 confessor, surnamed Trichinas, from the rough garment of hair which he wore. 
 He was renowned for many miracles, but especially for his power over the 
 demons. From his body issues a liquid that imparts health to the sick.
\switchcolumn*
\selectlanguage{latin}
In Monte Politiáno, in 
 Túscia, sanctæ Agnétis Vírginis, ex Ordine sancti Domínici, miráculis claræ.
\switchcolumn
\selectlanguage{english}
At Monte Pulciano, St. Agnes, a 
 virgin of the Order of St. Dominic, celebrated for her miracles.
\switchcolumn*
\selectlanguage{latin}
\end{paracol}


% ---- martyrology/mart04/mart0421.htm
\needspace{10\baselineskip}
\begin{paracol}{2}
\selectlanguage{latin}
\begin{center}{\color{gregoriocolor} Undécimo Kaléndas Maji. 
 Luna\dots\ }\end{center}
\switchcolumn
\selectlanguage{english}
\begin{center}{\color{gregoriocolor} The 
 21\textsuperscript{st} Day of April. The\dots\ Day of the Moon.}\end{center}
\end{paracol}

\noindent\begin{tabularx}{\linewidth}{*{19}{>{\centering\arraybackslash}X}}
 \textcolor{gregoriocolor}{a} & \textcolor{gregoriocolor}{b} & \textcolor{gregoriocolor}{c} & \textcolor{gregoriocolor}{d} & \textcolor{gregoriocolor}{e} & \textcolor{gregoriocolor}{f} & \textcolor{gregoriocolor}{g} & \textcolor{gregoriocolor}{h} & \textcolor{gregoriocolor}{i} & \textcolor{gregoriocolor}{k} & \textcolor{gregoriocolor}{l} & \textcolor{gregoriocolor}{m} & \textcolor{gregoriocolor}{n} & \textcolor{gregoriocolor}{p} & \textcolor{gregoriocolor}{q} & \textcolor{gregoriocolor}{r} & \textcolor{gregoriocolor}{s} & \textcolor{gregoriocolor}{t} & \textcolor{gregoriocolor}{u} \\
 23 & 24 & 25 & 26 & 27 & 28 & 29 & 1 & 2 & 3 & 4 & 5 & 6 & 7 & 8 & 9 & 10 & 11 & 12 \\
\end{tabularx}
\vspace{0.5\baselineskip}
\noindent\begin{tabularx}{\linewidth}{*{12}{>{\centering\arraybackslash}X}}
 \textcolor{gregoriocolor}{A} & \textcolor{gregoriocolor}{B} & \textcolor{gregoriocolor}{C} & \textcolor{gregoriocolor}{D} & \textcolor{gregoriocolor}{E} & F & \textcolor{gregoriocolor}{F} & \textcolor{gregoriocolor}{G} & \textcolor{gregoriocolor}{H} & \textcolor{gregoriocolor}{M} & \textcolor{gregoriocolor}{N} & \textcolor{gregoriocolor}{P} \\
 13 & 14 & 15 & 16 & 17 & 18 & 17 & 18 & 19 & 20 & 21 & 22 \\
\end{tabularx}

\begin{paracol}{2}
\selectlanguage{latin}
\lettrine[lines=2]{C}{antuáriæ,} in Anglia, 
 sancti Ansélmi Epíscopi, Confessóris et Ecclésiæ Doctóris, sanctitáte et 
 doctrína conspícui.
\switchcolumn
\selectlanguage{english}
\lettrine[lines=2]{A}{t} Canterbury, England, St. Anselm, 
 bishop, confessor, and doctor of the Church, renowned for sanctity and 
 learning.
\switchcolumn*
\selectlanguage{latin}
In Pérside natális 
 sancti Simeónis, Epíscopi Seleucíæ et Ctesiphóntis, qui, jubénte Rege 
 Persárum Sápore, comprehénsus ferróque onústus, iníquis tribunálibus 
 exhíbitus, et, cum Solem ipsum adoráre nollet et de Jesu Christo voce líbera 
 et constantíssima testarétur, primum carceráli ergástulo, cum áliis centum 
 (ex quibus álii Epíscopi, álii erant Presbyteri, álii diversórum órdinum 
 Clérici), longo témpore macerátus est. Deínde, cum Usthazánes, Regis 
 nutrítius, qui ante jam lapsus a fide, sed per eum ad pæniténtiam fúerat 
 revocátus, martyrium constánter subiísset, postrídie, qui erat ánnuus 
 Domínicæ passiónis dies, ómnibus ante Simeónis óculos, qui unumquémque eórum 
 strénue exhortabátur, gládio jugulátis, novíssime et ipse decollátus est. 
 Passi sunt étiam cum ipso claríssimi viri Abdéchalas et Ananías, qui ejus 
 erant Presbyteri, Pusícius quoque, Præféctus artíficum Regis, eo quod 
 Ananíam titubántem corroborásset, ídeo, collo circa téndinem perforáto et 
 lingua exínde extrácta, crudéli morte occúbuit; post quem et fília 
 cruciátibus ac demum ense decollátus est.
\switchcolumn
\selectlanguage{english}
In Persia, the birthday of St. 
 Simeon, bishop of Seleucia and Ctesiphon. He was arrested by order of 
 Sapor, king of Persia, loaded with irons, and presented to the iniquitous 
 tribunals. As he refused to adore the sun, and openly and constantly 
 bore testimony to Jesus Christ, he was confined for a long time in a dungeon 
 with one hundred other confessors, some of whom were bishops. others 
 priests, others clerics of various ranks. Afterwards, Usthazanes, the 
 king's foster-father, who had been converted from apostasy by Simeon, 
 endured martyrdom with great constancy. The day after, which was the 
 anniversary of our Lord's Passion, the companions of Simeon whom he had 
 feelingly exhorted, were beheaded before his eyes, after which he met the 
 same fate. With him suffered also several distinguished men: 
 Abdechalas and Ananias, his priests, with Pusicius, the head of the royal 
 workmen. This last having encouraged Ananias, who seemed to falter, 
 died a cruel death, having his tongue drawn out through a perforation made 
 in his neck. After him, his daughter, who was a consecrated virgin, 
 was put to death.
\switchcolumn*
\selectlanguage{latin}
Alexandríæ sanctórum 
 Mártyrum Aratóris Presbyteri, Fortunáti, Felícis, Sílvii et Vitális, qui in cárcere quievérunt.
\switchcolumn
\selectlanguage{english}
At Alexandria, the holy martyrs 
 Arator, a priest, Fortunatus, Felix, Silvius, and Vitalis, who all died in 
 prison.
\switchcolumn*
\selectlanguage{latin}
Nicomedíæ sanctórum 
 Mártyrum Apóllinis, Isácii et Codráti; e quibus, sub Diocletiáno Imperatóre, 
 últimus cápite plexus, et, paucis post illum diébus, duo primi in vínculis 
 fame confécti, martyrii corónam meruérunt.
\switchcolumn
\selectlanguage{english}
At Nicomedia, the holy martyrs 
 Apollo, Isacius, and Codratus, who suffered under the Emperor Diocletian. 
 The last of these was slain by the sword, and a few days later the other two 
 died from starvation in prison, meriting also the crown of martyrdom.
\switchcolumn*
\selectlanguage{latin}
Antiochíæ sancti 
 Anastásii Sinaítæ Epíscopi.
\switchcolumn
\selectlanguage{english}
At Antioch, St. Anastasius the 
 Sinaite, bishop.
\switchcolumn*
\selectlanguage{latin}
Œttingæ Véteris, in 
 Bavária, sancti Conrádi a Parzham, Confessóris, Ordinis Minórum Capuccinórum, 
 caritáte et oratióne insígnis; quem, miráculis clarum, Pius Papa Undécimus 
 Sanctórum número adscrípsit.
\switchcolumn
\selectlanguage{english}
At Wertingen in Bavaria, St. Conrad 
 of Parzham, confessor, of the Order of Friars Minor Capuchin, outstanding 
 both for prayer and for love of neighbour. Being renowned for 
 miracles, Pope Pius XI enrolled him among the number of the saints.
\switchcolumn*
\selectlanguage{latin}
\end{paracol}


% ---- martyrology/mart04/mart0422.htm
\needspace{10\baselineskip}
\begin{paracol}{2}
\selectlanguage{latin}
\begin{center}{\color{gregoriocolor} Décimo Kaléndas Maji. 
 Luna\dots\ }\end{center}
\switchcolumn
\selectlanguage{english}
\begin{center}{\color{gregoriocolor} The 
 22\textsuperscript{nd} Day of April. The\dots\ Day of the Moon.}\end{center}
\end{paracol}

\noindent\begin{tabularx}{\linewidth}{*{19}{>{\centering\arraybackslash}X}}
 \textcolor{gregoriocolor}{a} & \textcolor{gregoriocolor}{b} & \textcolor{gregoriocolor}{c} & \textcolor{gregoriocolor}{d} & \textcolor{gregoriocolor}{e} & \textcolor{gregoriocolor}{f} & \textcolor{gregoriocolor}{g} & \textcolor{gregoriocolor}{h} & \textcolor{gregoriocolor}{i} & \textcolor{gregoriocolor}{k} & \textcolor{gregoriocolor}{l} & \textcolor{gregoriocolor}{m} & \textcolor{gregoriocolor}{n} & \textcolor{gregoriocolor}{p} & \textcolor{gregoriocolor}{q} & \textcolor{gregoriocolor}{r} & \textcolor{gregoriocolor}{s} & \textcolor{gregoriocolor}{t} & \textcolor{gregoriocolor}{u} \\
 24 & 25 & 26 & 27 & 28 & 29 & 1 & 2 & 3 & 4 & 5 & 6 & 7 & 8 & 9 & 10 & 11 & 12 & 13 \\
\end{tabularx}
\vspace{0.5\baselineskip}
\noindent\begin{tabularx}{\linewidth}{*{12}{>{\centering\arraybackslash}X}}
 \textcolor{gregoriocolor}{A} & \textcolor{gregoriocolor}{B} & \textcolor{gregoriocolor}{C} & \textcolor{gregoriocolor}{D} & \textcolor{gregoriocolor}{E} & F & \textcolor{gregoriocolor}{F} & \textcolor{gregoriocolor}{G} & \textcolor{gregoriocolor}{H} & \textcolor{gregoriocolor}{M} & \textcolor{gregoriocolor}{N} & \textcolor{gregoriocolor}{P} \\
 14 & 15 & 16 & 17 & 18 & 19 & 18 & 19 & 20 & 21 & 22 & 23 \\
\end{tabularx}

\begin{paracol}{2}
\selectlanguage{latin}
\lettrine[lines=2]{R}{omæ,} via Appia, 
 natális sancti Sotéris, Papæ et Mártyris.
\switchcolumn
\selectlanguage{english}
\lettrine[lines=2]{A}{t} Rome, on the Appian Way, the 
 birthday of St. Soter, pope and martyr.
\switchcolumn*
\selectlanguage{latin}
Item Romæ sancti Caji, 
 Papæ et Mártyris; qui martyrio coronátus est sub Diocletiáno Príncipe.
\switchcolumn
\selectlanguage{english}
In the same city, Pope St. Caius, 
 who was crowned with martyrdom under Emperor Diocletian.
\switchcolumn*
\selectlanguage{latin}
Smyrnæ sanctórum 
 Apéllis et Lúcii, ex primis Christi discípulis.
\switchcolumn
\selectlanguage{english}
At Smyrna, the Saints Apelles and 
 Lucius, who were among the first disciples of Christ.
\switchcolumn*
\selectlanguage{latin}
Eódem die sanctórum 
 plurimórum Mártyrum, qui, sequénti anno post óbitum Simeónis, ánnuo item die 
 quo passiónis Domínicæ memória celebrabátur, per totam Pérsidis regiónem, 
 pro Christi nómine, sub Rege Sápore, gládio cædi jussi sunt. In quo 
 fídei certámine passus est Azades eunúchus, Regi caríssimus; Milles 
 Epíscopus, sanctitáte et miraculórum virtúte insígnis; Acépsimas Epíscopus, 
 cum Presbytero suo Jacóbo, item Aíthala et Josépho Presbyteris, Azadáne et 
 Abdiéso Diáconis, et complúribus áliis Cléricis; Maréas quoque et Bicor 
 Epíscopi, cum áliis vigínti Epíscopis, et Cléricis fere ducéntis 
 quinquagínta, Mónachis étiam et sacris Virgínibus plúrimis. Has inter 
 Vírgines fuit étiam sancti Simeónis Epíscopi soror, nómine Tárbula, cum 
 pedíssequa sua; quæ, stipítibus alligátæ serráque scissæ, crudelíssime 
 necátæ sunt.
\switchcolumn
\selectlanguage{english}
The same day, many holy martyrs 
 who, the year following the death of St. Simeon, and on the anniversary of 
 the Passion of our Lord, were put to the sword for the name of Christ 
 throughout Persia, under King Sapor. Among those who then suffered for 
 the faith were the eunuch Azades, a favorite of the king; Milles, a bishop 
 renowned for sanctity and miracles; Bishop Acepsimas with one of his priests 
 named James; also Aithalas and Joseph, priests; Azadan and Abdiesus, 
 deacons, and many other clerics; Mareas and Bicor, bishop, with twenty other 
 bishops, and nearly two hundred and fifty clerics; many monks and 
 consecrated virgins, among whom was the sister of St. Simeon, called Tarbula, 
 with her maid, who were both killed in a most cruel manner by being tied to 
 stakes and sawn asunder.
\switchcolumn*
\selectlanguage{latin}
Item in Pérside 
 sanctórum Parménii, Heliménæ et Chrysóteli Presbyterórum, Lucæ et Múcii 
 Diaconórum; quorum triúmphus martyrii in passióne sanctórum Abdon et Sennen 
 habétur.
\switchcolumn
\selectlanguage{english}
Also in Persia, Saints Parmenius, 
 Helimenas, and Chrysotelus, priests; Lucas and Mucius, deacons, whose 
 triumph is related in the Acts of Saints Abdon and Sennen.
\switchcolumn*
\selectlanguage{latin}
Alexandríæ natális 
 sancti Leónidæ Mártyris, qui sub Sevéro passus est.
\switchcolumn
\selectlanguage{english}
At Alexandria, the birthday of the 
 martyr St. Leonides, who suffered under Severus.
\switchcolumn*
\selectlanguage{latin}
Lugdúni, in Gállia, 
 sancti Epipódii, qui, in persecutióne Antoníni Veri, cum Alexándro colléga 
 tentus, ibídem, post dira torménta, martyrium abscissióne cápitis complévit.
\switchcolumn
\selectlanguage{english}
At Lyons in France, in the 
 persecution of Antoninus Verus, St. Epipodius, who was arrested with his 
 companion Alexander, and after undergoing severe torments, completed his 
 martyrdom by being beheaded.
\switchcolumn*
\selectlanguage{latin}
Constantinópoli sancti 
 Agapíti Papæ Primi, cujus sánctitas a beáto Gregório Magno commendátur. 
 Ipsíus autem corpus, póstea Romam relátum, in Vaticáno cónditum est.
\switchcolumn
\selectlanguage{english}
At Constantinople, Pope St. 
 Agapitus the First, whose sanctity was praised by St. Gregory the Great. 
 His body was afterwards taken to Rome and buried in the Vatican.
\switchcolumn*
\selectlanguage{latin}
Apud Senónas sancti 
 Leónis, Epíscopi et Confessóris.
\switchcolumn
\selectlanguage{english}
At Sens, St. Leo, bishop and 
 confessor.
\switchcolumn*
\selectlanguage{latin}
Anastasiópoli, in 
 Galátia, sancti Theodóri Epíscopi, miráculis clari.
\switchcolumn
\selectlanguage{english}
At Anastasiopolis in Galatia, St. 
 Theodore, a bishop well known for his miracles.
\switchcolumn*
\selectlanguage{latin}
\end{paracol}


% ---- martyrology/mart04/mart0423.htm
\needspace{10\baselineskip}
\begin{paracol}{2}
\selectlanguage{latin}
\begin{center}{\color{gregoriocolor} Nono Kaléndas Maji. 
 Luna\dots\ }\end{center}
\switchcolumn
\selectlanguage{english}
\begin{center}{\color{gregoriocolor} The 
 23\textsuperscript{rd} Day of April. The\dots\ Day of the Moon.}\end{center}
\end{paracol}

\noindent\begin{tabularx}{\linewidth}{*{19}{>{\centering\arraybackslash}X}}
 \textcolor{gregoriocolor}{a} & \textcolor{gregoriocolor}{b} & \textcolor{gregoriocolor}{c} & \textcolor{gregoriocolor}{d} & \textcolor{gregoriocolor}{e} & \textcolor{gregoriocolor}{f} & \textcolor{gregoriocolor}{g} & \textcolor{gregoriocolor}{h} & \textcolor{gregoriocolor}{i} & \textcolor{gregoriocolor}{k} & \textcolor{gregoriocolor}{l} & \textcolor{gregoriocolor}{m} & \textcolor{gregoriocolor}{n} & \textcolor{gregoriocolor}{p} & \textcolor{gregoriocolor}{q} & \textcolor{gregoriocolor}{r} & \textcolor{gregoriocolor}{s} & \textcolor{gregoriocolor}{t} & \textcolor{gregoriocolor}{u} \\
 25 & 26 & 27 & 28 & 29 & 1 & 2 & 3 & 4 & 5 & 6 & 7 & 8 & 9 & 10 & 11 & 12 & 13 & 14 \\
\end{tabularx}
\vspace{0.5\baselineskip}
\noindent\begin{tabularx}{\linewidth}{*{12}{>{\centering\arraybackslash}X}}
 \textcolor{gregoriocolor}{A} & \textcolor{gregoriocolor}{B} & \textcolor{gregoriocolor}{C} & \textcolor{gregoriocolor}{D} & \textcolor{gregoriocolor}{E} & F & \textcolor{gregoriocolor}{F} & \textcolor{gregoriocolor}{G} & \textcolor{gregoriocolor}{H} & \textcolor{gregoriocolor}{M} & \textcolor{gregoriocolor}{N} & \textcolor{gregoriocolor}{P} \\
 15 & 16 & 17 & 18 & 19 & 20 & 19 & 20 & 21 & 22 & 23 & 24 \\
\end{tabularx}

\begin{paracol}{2}
\selectlanguage{latin}
\lettrine[lines=2]{N}{atális} sancti Geórgii 
 Mártyris, cujus illústre martyrium inter Mártyrum corónas Ecclésia Dei 
 venerátur.
\switchcolumn
\selectlanguage{english}
\lettrine[lines=2]{T}{he} birthday of St. George, whose 
 illustrious martyrdom is honoured by the Church of God among the triumphs of 
 the other martyrs.
\switchcolumn*
\selectlanguage{latin}
In vico Tenkítten, ad 
 sinum Venédicum, in Borússia, item natális sancti Adalbérti, Epíscopi 
 Pragénsis et Mártyris, qui Polónis et Húngaris Evangélium prædicávit.
\switchcolumn
\selectlanguage{english}
At Danzig in Prussia, the birthday 
 of St. Adalbert, bishop of Prague, and martyr, who preached the Gospel to 
 the Poles and the Hungarians.
\switchcolumn*
\selectlanguage{latin}
Valéntiæ, in Gállia, 
 pássio sanctórum Mártyrum Felícis Presbyteri, Fortunáti et Achíllei 
 Diaconórum. Hi, cum fuíssent a beáto Irenæo, Lugdunénsi Episcopo, 
 missi ad prædicándum verbum Dei, et máximam illíus civitátis partem ad 
 Christi fidem convertíssent, a Duce Cornélio sunt in cárcerem trusi; deínde, 
 diutíssime verberáti, cruribúsque confráctis, circa rotárum vertíginem 
 stricti, fumum quoque in equúlei suspensióne perpéssi; ad extrémum gládio 
 consummáti sunt.
\switchcolumn
\selectlanguage{english}
At Valence in France, the holy 
 martyrs Felix, a priest, Fortunatus and Achilleus, deacons, who were sent 
 there to preach the word of God by blessed Irenæus, bishop of Lyons. 
 They converted the greater portion of that city to the faith of Christ. 
 These martyrs were cast into prison by the commander Cornelius, were for a 
 long time scourged, had their legs crushed, were bound to wheels in motion, 
 and stifled with smoke while stretched on the rack, and finally died by the 
 sword.
\switchcolumn*
\selectlanguage{latin}
Medioláni sancti 
 Mároli, Epíscopi et Confessóris.
\switchcolumn
\selectlanguage{english}
At Milan, St. Marolus, bishop and 
 confessor.
\switchcolumn*
\selectlanguage{latin}
Tulli, in Gállia, 
 sancti Gerárdi, ejúsdem civitátis Epíscopi.
\switchcolumn
\selectlanguage{english}
At Toul in France, St. Gerard, 
 bishop of that city.
\switchcolumn*
\selectlanguage{latin}
\end{paracol}


% ---- martyrology/mart04/mart0424.htm
\needspace{10\baselineskip}
\begin{paracol}{2}
\selectlanguage{latin}
\begin{center}{\color{gregoriocolor} Octávo Kaléndas Maji. 
 Luna\dots\ }\end{center}
\switchcolumn
\selectlanguage{english}
\begin{center}{\color{gregoriocolor} The 24\textsuperscript{th} Day of April. The\dots\ 
 Day of the Moon.}\end{center}
\end{paracol}

\noindent\begin{tabularx}{\linewidth}{*{19}{>{\centering\arraybackslash}X}}
 \textcolor{gregoriocolor}{a} & \textcolor{gregoriocolor}{b} & \textcolor{gregoriocolor}{c} & \textcolor{gregoriocolor}{d} & \textcolor{gregoriocolor}{e} & \textcolor{gregoriocolor}{f} & \textcolor{gregoriocolor}{g} & \textcolor{gregoriocolor}{h} & \textcolor{gregoriocolor}{i} & \textcolor{gregoriocolor}{k} & \textcolor{gregoriocolor}{l} & \textcolor{gregoriocolor}{m} & \textcolor{gregoriocolor}{n} & \textcolor{gregoriocolor}{p} & \textcolor{gregoriocolor}{q} & \textcolor{gregoriocolor}{r} & \textcolor{gregoriocolor}{s} & \textcolor{gregoriocolor}{t} & \textcolor{gregoriocolor}{u} \\
 26 & 27 & 28 & 29 & 1 & 2 & 3 & 4 & 5 & 6 & 7 & 8 & 9 & 10 & 11 & 12 & 13 & 14 & 15 \\
\end{tabularx}
\vspace{0.5\baselineskip}
\noindent\begin{tabularx}{\linewidth}{*{12}{>{\centering\arraybackslash}X}}
 \textcolor{gregoriocolor}{A} & \textcolor{gregoriocolor}{B} & \textcolor{gregoriocolor}{C} & \textcolor{gregoriocolor}{D} & \textcolor{gregoriocolor}{E} & F & \textcolor{gregoriocolor}{F} & \textcolor{gregoriocolor}{G} & \textcolor{gregoriocolor}{H} & \textcolor{gregoriocolor}{M} & \textcolor{gregoriocolor}{N} & \textcolor{gregoriocolor}{P} \\
 16 & 17 & 18 & 19 & 20 & 21 & 20 & 21 & 22 & 23 & 24 & 25 \\
\end{tabularx}

\begin{paracol}{2}
\selectlanguage{latin}
\lettrine[lines=2]{S}{evísii,} in Rhǽtia, 
 sancti Fidélis a Sigmarínga, Sacerdótis ex Ordine Minórum Capuccinórum et 
 Mártyris; qui, illuc ad prædicándam cathólicam fidem missus, ibídem, ab 
 hæréticis interémptus, martyrium consummávit; et a Benedícto Décimo quarto, 
 Pontifice Máximo, inter sanctos Mártyres relátus est.
\switchcolumn
\selectlanguage{english}
\lettrine[lines=2]{A}{t} Gruch in Switzerland, St. 
 Fidelis of Sigmaringen, priest and martyr, of the Order of Friars Minor 
 Capuchin. He was sent there to preach the Catholic faith, but was put 
 to death by the heretics. He was numbered among the holy martyrs by 
 the Sovereign Pontiff, Benedict XIV.
\switchcolumn*
\selectlanguage{latin}
Romæ sancti Sabæ, 
 ductóris mílitum, qui, accusátus quod Christiános in cárcere deténtos 
 visitáret, coram Júdice Christum líbere conféssus est. Hinc ab eódem 
 Júdice fácibus adústus et in lebétem picis fervéntis est immíssus, et, cum 
 inde evasísset illæsus, eo miráculo septuagínta viros ad Christum convértit; 
 qui omnes, constánter in confessióne fídei permanéntes, gládio cæsi sunt. 
 Postrémo et ipse, demérsus in flumen, martyrium consummávit.
\switchcolumn
\selectlanguage{english}
At Rome, St. Sabas, a military 
 officer, who bravely confessed Christ before the judge when he was accused 
 of visiting the Christians kept in prison. For this he was burned with 
 torches and thrown into a cauldron of boiling pitch, out of which he came 
 uninjured. Seventy men were converted to Christ at the sight of this 
 miracle, and as they all remained unshaken in the confession of the faith, 
 they were put to the sword. Sabas, however, completed his martyrdom by 
 being cast into the river.
\switchcolumn*
\selectlanguage{latin}
Lugdúni, in Gállia, 
 natális sancti Alexándri Mártyris, qui, in persecutióne Antoníni Veri, post cárceris custódiam, primo ita laniátus est crudelitáte verberántium, ut, 
 crate solúta costárum, patefáctis viscéribus, interióra córporis panderéntur; 
 deínde, crucis affixus patíbulo, beátum spíritum exanimátus emísit. 
 Passi sunt cum ipso et álii, número trigínta quátuor, quorum memória áliis 
 diébus ágitur.
\switchcolumn
\selectlanguage{english}
At Lyons in France, during the 
 persecution of Antoninus Verus, the birthday of St. Alexander, martyr. 
 After being imprisoned, he was so lacerated by the cruelty of those who 
 scourged him, that his ribs and the interior of his body were exposed to 
 view. Then he was fastened to the gibbet of the cross, on which he 
 yielded up his blessed soul. Thirty-four others who suffered with him 
 are commemorated on other days.
\switchcolumn*
\selectlanguage{latin}
Nicomedíæ sanctórum 
 Mártyrum Eusébii, Neónis, Leóntii, Longíni et aliórum quátuor; qui, in 
 persecutióne Diocletiáni, post diros cruciátus, gládio percússi sunt.
\switchcolumn
\selectlanguage{english}
At Nicomedia, during the 
 persecution of Diocletian, the holy martyrs Eusebius, Neon, Leontius, 
 Longinus, and four others, all of whom were slain with the sword after 
 enduring great torments.
\switchcolumn*
\selectlanguage{latin}
In Anglia deposítio 
 sancti Mellíti Epíscopi, qui, a sancto Gregório Papa in Angliam missus, 
 orientáles Saxónes et ipsórum Regem ad fidem convértit.
\switchcolumn
\selectlanguage{english}
In England, the death of St. 
 Mellitus, bishop. He was sent there by St. Gregory, and he converted 
 to the faith the East Saxons and their king.
\switchcolumn*
\selectlanguage{latin}
Illíberi, in Hispánia, sancti Gregórii, Epíscopi et Confessóris.
\switchcolumn
\selectlanguage{english}
At Elvira, in Spain, St. Gregory, 
 bishop and confessor.
\switchcolumn*
\selectlanguage{latin}
Bríxiæ sancti Honórii 
 Epíscopi.
\switchcolumn
\selectlanguage{english}
At Brescia, St. Honorius, bishop.
\switchcolumn*
\selectlanguage{latin}
In Ióna, Scótiæ 
 ínsula, 
 sancti Egbérti, Presbyteri et Mónachi, admirándæ humilitátis et continéntiæ 
 viri.
\switchcolumn
\selectlanguage{english}
In Iona, an island of Scotland, St. 
 Egbert, priest and monk, a man of admirable humility and continency.
\switchcolumn*
\selectlanguage{latin}
Rhemis, in Gállia, 
 sanctárum Vírginum Bovæ et Dodæ.
\switchcolumn
\selectlanguage{english}
At Rheims in France, the holy 
 virgins Bova and Doda.
\switchcolumn*
\selectlanguage{latin}
Andégavi, in Gállia, 
 sanctæ Maríæ a sancta Euphrásia Pelletier, Vírginis, Institúti Sorórum a 
 Bono Pastóre Fundatrícis, quam Pius Duodécimus, Póntifex Máximus, in 
 Sanctárum númerum rétulit.
\switchcolumn
\selectlanguage{english}
At Angers in France, St. Mary Euphrasia Pelletier, virgin and foundress of the Institute of the Good 
 Shepherd Sisters, whom Pius XII, Sovereign Pontiff, enrolled among the 
 number of the saints.
\switchcolumn*
\selectlanguage{latin}
Medioláni Convérsio 
 sancti Augustíni Epíscopi, Confessóris et Ecclésiæ Doctóris; quem beátum 
 Ambrósius Epíscopus veritátem fídei cathólicæ dócuit, et hac die baptizávit.
\switchcolumn
\selectlanguage{english}
At Milan, the Conversion of St. 
 Augustine, bishop, confessor, and doctor of the Church, whom the bishop St. 
 Ambrose had instructed in the truth of the Catholic faith, and baptized on 
 this day.
\switchcolumn*
\selectlanguage{latin}
\end{paracol}


% ---- martyrology/mart04/mart0425.htm
\needspace{10\baselineskip}
\begin{paracol}{2}
\selectlanguage{latin}
\begin{center}{\color{gregoriocolor} Séptimo Kaléndas Maji. 
 Luna\dots\ }\end{center}
\switchcolumn
\selectlanguage{english}
\begin{center}{\color{gregoriocolor} The 
 25\textsuperscript{th} Day of April. The\dots\ Day of the Moon.}\end{center}
\end{paracol}

\noindent\begin{tabularx}{\linewidth}{*{19}{>{\centering\arraybackslash}X}}
 \textcolor{gregoriocolor}{a} & \textcolor{gregoriocolor}{b} & \textcolor{gregoriocolor}{c} & \textcolor{gregoriocolor}{d} & \textcolor{gregoriocolor}{e} & \textcolor{gregoriocolor}{f} & \textcolor{gregoriocolor}{g} & \textcolor{gregoriocolor}{h} & \textcolor{gregoriocolor}{i} & \textcolor{gregoriocolor}{k} & \textcolor{gregoriocolor}{l} & \textcolor{gregoriocolor}{m} & \textcolor{gregoriocolor}{n} & \textcolor{gregoriocolor}{p} & \textcolor{gregoriocolor}{q} & \textcolor{gregoriocolor}{r} & \textcolor{gregoriocolor}{s} & \textcolor{gregoriocolor}{t} & \textcolor{gregoriocolor}{u} \\
 27 & 28 & 29 & 1 & 2 & 3 & 4 & 5 & 6 & 7 & 8 & 9 & 10 & 11 & 12 & 13 & 14 & 15 & 16 \\
\end{tabularx}
\vspace{0.5\baselineskip}
\noindent\begin{tabularx}{\linewidth}{*{12}{>{\centering\arraybackslash}X}}
 \textcolor{gregoriocolor}{A} & \textcolor{gregoriocolor}{B} & \textcolor{gregoriocolor}{C} & \textcolor{gregoriocolor}{D} & \textcolor{gregoriocolor}{E} & F & \textcolor{gregoriocolor}{F} & \textcolor{gregoriocolor}{G} & \textcolor{gregoriocolor}{H} & \textcolor{gregoriocolor}{M} & \textcolor{gregoriocolor}{N} & \textcolor{gregoriocolor}{P} \\
 17 & 18 & 19 & 20 & 21 & 22 & 21 & 22 & 23 & 24 & 25 & 26 \\
\end{tabularx}

\begin{paracol}{2}
\selectlanguage{latin}
\lettrine[lines=1]{R}{omæ} Litaníæ majóres ad sanctum Petrum.
\switchcolumn
\selectlanguage{english}
\lettrine[lines=1]{A}{t} Rome, the Greater Litanies at St. Peter's.
\switchcolumn*
\selectlanguage{latin}
Alexandríæ natális 
 beáti Marci Evangelístæ. Hic, discípulus et intérpres Apóstoli Petri, 
 rogátus Romæ a frátribus scripsit Evangélium, quo assúmpto, perréxit in 
 Ægyptum, primúsque Alexandríæ Christum annúntians, constítuit Ecclésiam; ac 
 póstea, pro fide Christi tentus, fúnibus vinctus et per saxa raptátus, 
 gráviter afflíctus est; deínde, reclúsus in cárcere, primo angélica visitatióne confortátus est, et demum, ipso Dómino sibi apparénte, ad 
 cæléstia regna vocátus, octávo Nerónis anno.
\switchcolumn
\selectlanguage{english}
At Alexandria, the birthday of St. 
 Mark the Evangelist, disciple and interpreter of the apostle St. Peter. 
 He wrote his gospel at the request of the faithful at Rome, and taking it 
 with him, proceeded to Egypt and founded a church at Alexandria, where he 
 was the first to preach Christ. Afterwards, being arrested for the 
 faith, he was bound, dragged over stones, and endured great afflictions. 
 Finally he was confined to prison, where, being comforted by the visit of an 
 angel, and even by an apparition of our Lord himself, he was called to the 
 heavenly kingdom in the eighth year of the reign of Nero.
\switchcolumn*
\selectlanguage{latin}
Item Alexandríæ sancti 
 Aniáni Epíscopi, qui, beáti Marci discípulus ejúsque in Episcopátu 
 succéssor, clarus virtútibus quiévit in Dómino.
\switchcolumn
\selectlanguage{english}
Also at Alexandria, Bishop St. 
 Anian, disciple of blessed Mark, and his successor in the episcopate. 
 With a great renown for virtue, he rested in the Lord.
\switchcolumn*
\selectlanguage{latin}
Antiochíæ sancti 
 Stéphani, Epíscopi et Mártyris, qui ab hæréticis Synodum Chalcedonénsem 
 impugnántibus, multa passus, in Oróntem flúvium præcipitátus est, témpore 
 Zenónis Imperatóris.
\switchcolumn
\selectlanguage{english}
At Antioch, St. Stephen, bishop and 
 martyr, who suffered a great deal from the heretics opposed to the Council 
 of Chalcedon, and was cast into the river Orontes, in the time of Emperor 
 Zeno.
\switchcolumn*
\selectlanguage{latin}
Syracúsis, in Sicília, 
 sanctórum Mártyrum fratrum Evódii, Hermógenis et Callístæ.
\switchcolumn
\selectlanguage{english}
At Syracuse in Sicily, the holy 
 martyrs Evodius, Hermogenes, and Callista.
\switchcolumn*
\selectlanguage{latin}
Láubiis, in Bélgio, 
 natális sancti Ermíni, Epíscopi et Confessóris.
\switchcolumn
\selectlanguage{english}
At Lobbes in Belgium, the birthday 
 of St. Ermin, bishop and confessor.
\switchcolumn*
\selectlanguage{latin}
Antiochíæ sanctórum 
 Philónis et Agathópodis Diaconórum, de quibus beátus Ignátius, Epíscopus et 
 Martyr, laudábilem in suis epístolis mentiónem facit.
\switchcolumn
\selectlanguage{english}
At Antioch, the deacons Saints 
 Philo and Agathopodes, who were praised in the letters of blessed Ignatius, 
 bishop and martyr.
\switchcolumn*
\selectlanguage{latin}
\end{paracol}


% ---- martyrology/mart04/mart0426.htm
\needspace{10\baselineskip}
\begin{paracol}{2}
\selectlanguage{latin}
\begin{center}{\color{gregoriocolor} Sexto Kaléndas Maji. 
 Luna\dots\ }\end{center}
\switchcolumn
\selectlanguage{english}
\begin{center}{\color{gregoriocolor} The 
 26\textsuperscript{th} Day of April. The\dots\ Day of the Moon.}\end{center}
\end{paracol}

\noindent\begin{tabularx}{\linewidth}{*{19}{>{\centering\arraybackslash}X}}
 \textcolor{gregoriocolor}{a} & \textcolor{gregoriocolor}{b} & \textcolor{gregoriocolor}{c} & \textcolor{gregoriocolor}{d} & \textcolor{gregoriocolor}{e} & \textcolor{gregoriocolor}{f} & \textcolor{gregoriocolor}{g} & \textcolor{gregoriocolor}{h} & \textcolor{gregoriocolor}{i} & \textcolor{gregoriocolor}{k} & \textcolor{gregoriocolor}{l} & \textcolor{gregoriocolor}{m} & \textcolor{gregoriocolor}{n} & \textcolor{gregoriocolor}{p} & \textcolor{gregoriocolor}{q} & \textcolor{gregoriocolor}{r} & \textcolor{gregoriocolor}{s} & \textcolor{gregoriocolor}{t} & \textcolor{gregoriocolor}{u} \\
 28 & 29 & 1 & 2 & 3 & 4 & 5 & 6 & 7 & 8 & 9 & 10 & 11 & 12 & 13 & 14 & 15 & 16 & 17 \\
\end{tabularx}
\vspace{0.5\baselineskip}
\noindent\begin{tabularx}{\linewidth}{*{12}{>{\centering\arraybackslash}X}}
 \textcolor{gregoriocolor}{A} & \textcolor{gregoriocolor}{B} & \textcolor{gregoriocolor}{C} & \textcolor{gregoriocolor}{D} & \textcolor{gregoriocolor}{E} & F & \textcolor{gregoriocolor}{F} & \textcolor{gregoriocolor}{G} & \textcolor{gregoriocolor}{H} & \textcolor{gregoriocolor}{M} & \textcolor{gregoriocolor}{N} & \textcolor{gregoriocolor}{P} \\
 18 & 19 & 20 & 21 & 22 & 23 & 22 & 23 & 24 & 25 & 26 & 27 \\
\end{tabularx}

\begin{paracol}{2}
\selectlanguage{latin}
\lettrine[lines=2]{R}{omæ} natális beáti 
 Cleti, Papæ et Mártyris; qui, secúndus post Apóstolum Petrum, rexit 
 Ecclésiam, et martyrio in persecutióne Domitiáni coronátus est.
\switchcolumn
\selectlanguage{english}
\lettrine[lines=2]{A}{t} Rome, the birthday of St. 
 Cletus, the pope who governed the Church the second after the apostle St. 
 Peter, and was crowned with martyrdom in the persecution of Domitian.
\switchcolumn*
\selectlanguage{latin}
Sancti Marcellíni, 
 Papæ et Mártyris, cujus dies natális octávo Kaléndas Novémbris recensétur.
\switchcolumn
\selectlanguage{english}
St. Marcellinus, pope and martyr, 
 whose birthday is commemorated on the 25th of October.
\switchcolumn*
\selectlanguage{latin}
Amaséæ, in Ponto, 
 sancti Basiléi, Epíscopi et Mártyris, qui, sub Licínio Imperatóre, illústre 
 martyrium consummávit. Ipsíus autem corpus, in mare projéctum, et ab 
 Elpidíphoro, Angeli mónitu repértum, honorífice tumulátum fuit.
\switchcolumn
\selectlanguage{english}
At Amasea in Pontus, St. Basileus, 
 bishop and martyr, whose illustrious martyrdom occurred under Emperor 
 Licinius. His body was thrown into the sea, but was found by 
 Elpidiphorus, through the revelation of an angel, and was honorably buried.
\switchcolumn*
\selectlanguage{latin}
Brácari, in Lusitánia, 
 sancti Petri Mártyris, qui fuit primus ejúsdem civitátis Epíscopus.
\switchcolumn
\selectlanguage{english}
At Braga in Portugal, St. Peter, 
 martyr, the first bishop of that city.
\switchcolumn*
\selectlanguage{latin}
Viénnæ, in Gállia, 
 sancti Claréntii, Epíscopi et Confessóris.
\switchcolumn
\selectlanguage{english}
At Vienne in France, St. Clarence, 
 bishop and confessor.
\switchcolumn*
\selectlanguage{latin}
Verónæ sancti Lucídii 
 Epíscopi.
\switchcolumn
\selectlanguage{english}
At Verona, St. Lucidius, bishop.
\switchcolumn*
\selectlanguage{latin}
In monastério Céntula, 
 in Gállia, sancti Richárii, Presbyteri et Confessóris.
\switchcolumn
\selectlanguage{english}
In the monastery of Centula in 
 France, St. Richarius, priest and confessor.
\switchcolumn*
\selectlanguage{latin}
Trecis, in Gállia, 
 sanctæ Exsuperántiæ Vírginis.
\switchcolumn
\selectlanguage{english}
At Troyes in France, St. 
 Exuperantia, virgin.
\switchcolumn*
\selectlanguage{latin}
\end{paracol}


% ---- martyrology/mart04/mart0427.htm
\needspace{10\baselineskip}
\begin{paracol}{2}
\selectlanguage{latin}
\begin{center}{\color{gregoriocolor} Quinto Kaléndas Maji. 
 Luna\dots\ }\end{center}
\switchcolumn
\selectlanguage{english}
\begin{center}{\color{gregoriocolor} The 
 27\textsuperscript{th} Day of April. The\dots\ Day of the Moon.}\end{center}
\end{paracol}

\noindent\begin{tabularx}{\linewidth}{*{19}{>{\centering\arraybackslash}X}}
 \textcolor{gregoriocolor}{a} & \textcolor{gregoriocolor}{b} & \textcolor{gregoriocolor}{c} & \textcolor{gregoriocolor}{d} & \textcolor{gregoriocolor}{e} & \textcolor{gregoriocolor}{f} & \textcolor{gregoriocolor}{g} & \textcolor{gregoriocolor}{h} & \textcolor{gregoriocolor}{i} & \textcolor{gregoriocolor}{k} & \textcolor{gregoriocolor}{l} & \textcolor{gregoriocolor}{m} & \textcolor{gregoriocolor}{n} & \textcolor{gregoriocolor}{p} & \textcolor{gregoriocolor}{q} & \textcolor{gregoriocolor}{r} & \textcolor{gregoriocolor}{s} & \textcolor{gregoriocolor}{t} & \textcolor{gregoriocolor}{u} \\
 29 & 1 & 2 & 3 & 4 & 5 & 6 & 7 & 8 & 9 & 10 & 11 & 12 & 13 & 14 & 15 & 16 & 17 & 18 \\
\end{tabularx}
\vspace{0.5\baselineskip}
\noindent\begin{tabularx}{\linewidth}{*{12}{>{\centering\arraybackslash}X}}
 \textcolor{gregoriocolor}{A} & \textcolor{gregoriocolor}{B} & \textcolor{gregoriocolor}{C} & \textcolor{gregoriocolor}{D} & \textcolor{gregoriocolor}{E} & F & \textcolor{gregoriocolor}{F} & \textcolor{gregoriocolor}{G} & \textcolor{gregoriocolor}{H} & \textcolor{gregoriocolor}{M} & \textcolor{gregoriocolor}{N} & \textcolor{gregoriocolor}{P} \\
 19 & 20 & 21 & 22 & 23 & 24 & 23 & 24 & 25 & 26 & 27 & 28 \\
\end{tabularx}

\begin{paracol}{2}
\selectlanguage{latin}
\lettrine[lines=2]{S}{ancti} Petri Canísii, 
 Sacerdótis e Societáte Jesu et Confessóris atque Ecclésiæ Doctóris; qui 
 duodécimo Kaléndas Januárii migrávit ad Dóminum.
\switchcolumn
\selectlanguage{english}
\lettrine[lines=2]{S}{t.} Peter Canisius, priest of the 
 Society of Jesus, confessor and doctor of the Church, who departed to the 
 Lord on the 21st of December.
\switchcolumn*
\selectlanguage{latin}
Nicomedíæ natális 
 sancti Anthimi, Epíscopi et Mártyris; qui in persecutióne Diocletiáni, ob 
 confessiónem Christi, martyrii glóriam obtruncatióne cápitis accépit. 
 Secúta est quoque illum univérsa ferme gregis sui multitúdo; quorum álios 
 Judex gládio obtruncári, álios conflagrári ígnibus, álios, navículis 
 impósitos, pélago immérgi fecit.
\switchcolumn
\selectlanguage{english}
At Nicomedia, during the 
 persecution of Diocletian, the birthday of St. Anthimus, bishop and martyr, 
 who obtained the glory of martyrdom by being beheaded for the faith. 
 Nearly all his numerous flock followed him. The judge ordered some to 
 be beheaded, others to be burned alive, others to be put in boats and sunk 
 in the sea.
\switchcolumn*
\selectlanguage{latin}
Tarsi, in Cilícia, 
 sanctórum Cástoris et Stéphani Mártyrum.
\switchcolumn
\selectlanguage{english}
At Tarsus in Cilicia, the Saints 
 Castor and Stephen, martyrs.
\switchcolumn*
\selectlanguage{latin}
Bonóniæ sancti 
 Tertulliáni, Epíscopi et Confessóris.
\switchcolumn
\selectlanguage{english}
At Bologna, St. Tertullian, bishop 
 and confessor.
\switchcolumn*
\selectlanguage{latin}
Bríxiæ sancti 
 Theóphili Epíscopi.
\switchcolumn
\selectlanguage{english}
At Brescia, St. Theophilus, bishop.
\switchcolumn*
\selectlanguage{latin}
In Ægypto sancti 
 Theodóri Abbátis, qui fuit discípulus sancti Pachómii.
\switchcolumn
\selectlanguage{english}
In Egypt, St. Theodore, abbot, who 
 was a disciple of St. Pachomius.
\switchcolumn*
\selectlanguage{latin}
Constantinópoli sancti 
 Joánnis Abbátis, qui pro cultu sacrárum Imáginum, sub Leóne Isáurico, 
 plúrimum decertávit.
\switchcolumn
\selectlanguage{english}
At Constantinople, the abbot St. 
 John, who valiantly defended the veneration of sacred images, under Leo the 
 Isaurian.
\switchcolumn*
\selectlanguage{latin}
Tarracóne, in 
 Hispánia, beáti Petri Armengáudii, ex Ordine beátæ Maríæ de Mercéde 
 redemptiónis captivórum; qui, multa pro fidélibus rediméndis in Africa 
 passus, tandem, in convéntu sanctæ Maríæ Pratórum, beáto fine quiévit.
\switchcolumn
\selectlanguage{english}
At Tarragona in Spain, the blessed 
 Peter Armengaudius, of the Order of Blessed Mary of Mercy for the Redemption 
 of Captives. He endured many tribulations in Africa in ransoming the 
 faithful, and finally closed his career peacefully in the convent of St. 
 Mary of the Meadows.
\switchcolumn*
\selectlanguage{latin}
Lucæ, in Túscia, beátæ 
 Zitæ Vírginis, virtútum et miraculórum fama conspícuæ.
\switchcolumn
\selectlanguage{english}
At Lucca in Tuscany, blessed Zita, a 
 virgin renowned for virtues and miracles.
\switchcolumn*
\selectlanguage{latin}
\end{paracol}


% ---- martyrology/mart04/mart0428.htm
\needspace{10\baselineskip}
\begin{paracol}{2}
\selectlanguage{latin}
\begin{center}{\color{gregoriocolor} Quarto Kaléndas Maji. 
 Luna\dots\ }\end{center}
\switchcolumn
\selectlanguage{english}
\begin{center}{\color{gregoriocolor} The 
 28\textsuperscript{th} Day of April. The\dots\ Day of the Moon.}\end{center}
\end{paracol}

\noindent\begin{tabularx}{\linewidth}{*{19}{>{\centering\arraybackslash}X}}
 \textcolor{gregoriocolor}{a} & \textcolor{gregoriocolor}{b} & \textcolor{gregoriocolor}{c} & \textcolor{gregoriocolor}{d} & \textcolor{gregoriocolor}{e} & \textcolor{gregoriocolor}{f} & \textcolor{gregoriocolor}{g} & \textcolor{gregoriocolor}{h} & \textcolor{gregoriocolor}{i} & \textcolor{gregoriocolor}{k} & \textcolor{gregoriocolor}{l} & \textcolor{gregoriocolor}{m} & \textcolor{gregoriocolor}{n} & \textcolor{gregoriocolor}{p} & \textcolor{gregoriocolor}{q} & \textcolor{gregoriocolor}{r} & \textcolor{gregoriocolor}{s} & \textcolor{gregoriocolor}{t} & \textcolor{gregoriocolor}{u} \\
 1 & 2 & 3 & 4 & 5 & 6 & 7 & 8 & 9 & 10 & 11 & 12 & 13 & 14 & 15 & 16 & 17 & 18 & 19 \\
\end{tabularx}
\vspace{0.5\baselineskip}
\noindent\begin{tabularx}{\linewidth}{*{12}{>{\centering\arraybackslash}X}}
 \textcolor{gregoriocolor}{A} & \textcolor{gregoriocolor}{B} & \textcolor{gregoriocolor}{C} & \textcolor{gregoriocolor}{D} & \textcolor{gregoriocolor}{E} & F & \textcolor{gregoriocolor}{F} & \textcolor{gregoriocolor}{G} & \textcolor{gregoriocolor}{H} & \textcolor{gregoriocolor}{M} & \textcolor{gregoriocolor}{N} & \textcolor{gregoriocolor}{P} \\
 20 & 21 & 22 & 23 & 24 & 25 & 24 & 25 & 26 & 27 & 28 & 29 \\
\end{tabularx}

\begin{paracol}{2}
\selectlanguage{latin}
\lettrine[lines=2]{S}{ancti} Pauli a Cruce, 
 Presbyteri et Confessóris; qui Congregatiónis a Cruce et Passióne Dómini 
 nostri Jesu Christi nuncupátæ Institútor fuit, atque in Dómino obdormívit 
 quintodécimo Kaléndas Novémbris.
\switchcolumn
\selectlanguage{english}
\lettrine[lines=2]{S}{t.} Paul of the Cross, priest and 
 confessor, founder of the Congregation of the Cross and Passion of our Lord 
 Jesus Christ. He went to his repose in the Lord on the 18th of 
 October.
\switchcolumn*
\selectlanguage{latin}
Ravénnæ natális sancti 
 Vitális Mártyris, viri sanctæ Valériæ ac patris sanctórum Gervásii et 
 Protásii; qui, cum beáti Ursicíni corpus sublátum débita honestáte 
 sepelísset, tentus est a Paulíno Consulári, et, post equúlei torménta, 
 jussus depóni in profúndam fóveam, et terra ac lapídibus óbrui; talíque 
 martyrio migrávit ad Christum.
\switchcolumn
\selectlanguage{english}
At Ravenna, the birthday of St. 
 Vitalis, martyr, father of the Saints Gervase and Protase. When he had 
 taken up and reverently buried the body of blessed Ursicinus, he was 
 arrested by the governor Paulinus, and after being racked and thrown into a 
 deep pit, was covered with earth and stones, and by this kind of martyrdom went 
 to Christ.
\switchcolumn*
\selectlanguage{latin}
Atínæ, in Campánia, 
 sancti Marci, qui, a beáto Petro Apóstolo Epíscopus ordinátus, Æquícolis 
 primus Evangélium prædicávit, et in persecutióne Domitiáni, sub Máximo 
 Præside, martyrii corónam accépit.
\switchcolumn
\selectlanguage{english}
At Atino in Campania, St. Mark, who 
 was made bishop by the blessed apostle Peter. He was the first to 
 preach the Gospel to the Equicoli, and received the crown of martyrdom in 
 the persecution of Domitian, under the governor Maximus.
\switchcolumn*
\selectlanguage{latin}
Prusæ, in Bithynia, 
 sanctórum Mártyrum Patrícii Epíscopi, Acátii, Menándri et Polyǽni.
\switchcolumn
\selectlanguage{english}
At Broussa in Bithynia, the holy 
 martyrs Patrick, a bishop, Acatius, Menander, and Polyaenus.
\switchcolumn*
\selectlanguage{latin}
Eódem die sanctórum 
 Mártyrum Aphrodísii, Caralíppi, Agápii et Eusébii.
\switchcolumn
\selectlanguage{english}
On the same day, the holy martyrs 
 Aphrodisius, Caralippus, Agapius, and Eusebius.
\switchcolumn*
\selectlanguage{latin}
In Pannónia sancti 
 Polliónis Mártyris, sub Diocletiáno Imperatóre.
\switchcolumn
\selectlanguage{english}
In Hungary, St. Pollio, martyr, 
 under the Emperor Diocletian.
\switchcolumn*
\selectlanguage{latin}
Medioláni sanctæ 
 Valériæ Mártyris, uxóris sancti Vitális ac matris sanctórum Gervásii et 
 Protásii.
\switchcolumn
\selectlanguage{english}
At Milan, the martyr St. Valeria, 
 who was the wife of St. Vitalis and the mother of Saints Gervase and Protase.
\switchcolumn*
\selectlanguage{latin}
Alexandríæ pássio 
 sanctæ Theodóræ, Vírginis et Mártyris. Hæc, idólis sacrificáre 
 contémnens, in lupánar est trádita, sed repénte quidam ex frátribus, nómine 
 Dídymus, miro Dei favóre, commutátis véstibus, illam erípuit; qui póstea, in 
 persecutióne Diocletiáni, sub Eustrátio Prǽside, simul cum eádem Vírgine 
 percússus, simul coronátus est.
\switchcolumn
\selectlanguage{english}
At Alexandria, the martyrdom of the 
 virgin St. Theodora. For refusing to sacrifice to idols, she was sent 
 to a place of debauchery; but one of the brethren, named Didymus, through 
 the admirable providence of God, delivered her by quickly exchanging 
 garments with her. He was afterwards beheaded and crowned with her in 
 the persecution of Diocletian, under the governor Eustratius.
\switchcolumn*
\selectlanguage{latin}
Turiasóne, in Hispánia 
 Tarraconénsi, sancti Prudéntii, Epíscopi et Confessóris.
\switchcolumn
\selectlanguage{english}
At Tarrazona in Spain, St. 
 Prudentius, bishop and confessor.
\switchcolumn*
\selectlanguage{latin}
Corfínii, in Pelígnis, 
 sancti Pámphili, Valvénsis Epíscopi, caritáte in páuperes et virtúte 
 miraculórum illústris; cujus corpus Sulmóne cónditum est.
\switchcolumn
\selectlanguage{english}
At Corfinio in Peligno, St. 
 Pamphilus, bishop of Valva, illustrious for his charity towards the poor and 
 the gift of miracles. His body was buried at Solmona.
\switchcolumn*
\selectlanguage{latin}
In pago sancti 
 Lauréntii ad Séparim, diœcésis Lucionénsis, sancti Ludovíci Maríæ Grignion a 
 Montfort, Confessóris, Fundatóris Missionariórum Societátis Maríæ et 
 Filiárum a Sapiéntia, apostólicæ vitæ forma, prædicatióne et devotióne 
 mariáli insígnis, quem Pius Papa Duodécimus Sanctórum catálogo adscrípsit.
\switchcolumn
\selectlanguage{english}
At St. Laurent sur Sèvres, in the 
 diocese of Luçon, St. Louis-Marie Grignion de Montfort, confessor and 
 founder of the Missionaries of the Company of Mary and the Sisters of 
 Wisdom, a form of apostolic life. He was renowned for his preaching 
 and devotion to the Blessed Mother, and was added to the number of the 
 saints by Pope Pius XII.
\switchcolumn*
\selectlanguage{latin}
\end{paracol}


% ---- martyrology/mart04/mart0429.htm
\needspace{10\baselineskip}
\begin{paracol}{2}
\selectlanguage{latin}
\begin{center}{\color{gregoriocolor} Tértio Kaléndas Maji. 
 Luna\dots\ }\end{center}
\switchcolumn
\selectlanguage{english}
\begin{center}{\color{gregoriocolor} The 
 29\textsuperscript{th} Day of April. The\dots\ Day of the Moon.}\end{center}
\end{paracol}

\noindent\begin{tabularx}{\linewidth}{*{19}{>{\centering\arraybackslash}X}}
 \textcolor{gregoriocolor}{a} & \textcolor{gregoriocolor}{b} & \textcolor{gregoriocolor}{c} & \textcolor{gregoriocolor}{d} & \textcolor{gregoriocolor}{e} & \textcolor{gregoriocolor}{f} & \textcolor{gregoriocolor}{g} & \textcolor{gregoriocolor}{h} & \textcolor{gregoriocolor}{i} & \textcolor{gregoriocolor}{k} & \textcolor{gregoriocolor}{l} & \textcolor{gregoriocolor}{m} & \textcolor{gregoriocolor}{n} & \textcolor{gregoriocolor}{p} & \textcolor{gregoriocolor}{q} & \textcolor{gregoriocolor}{r} & \textcolor{gregoriocolor}{s} & \textcolor{gregoriocolor}{t} & \textcolor{gregoriocolor}{u} \\
 2 & 3 & 4 & 5 & 6 & 7 & 8 & 9 & 10 & 11 & 12 & 13 & 14 & 15 & 16 & 17 & 18 & 19 & 20 \\
\end{tabularx}
\vspace{0.5\baselineskip}
\noindent\begin{tabularx}{\linewidth}{*{12}{>{\centering\arraybackslash}X}}
 \textcolor{gregoriocolor}{A} & \textcolor{gregoriocolor}{B} & \textcolor{gregoriocolor}{C} & \textcolor{gregoriocolor}{D} & \textcolor{gregoriocolor}{E} & F & \textcolor{gregoriocolor}{F} & \textcolor{gregoriocolor}{G} & \textcolor{gregoriocolor}{H} & \textcolor{gregoriocolor}{M} & \textcolor{gregoriocolor}{N} & \textcolor{gregoriocolor}{P} \\
 21 & 22 & 23 & 24 & 25 & 26 & 25 & 26 & 27 & 28 & 29 & 1 \\
\end{tabularx}

\begin{paracol}{2}
\selectlanguage{latin}
\lettrine[lines=2]{S}{ancti} Petri, ex 
 Ordine Prædicatórum, Mártyris, qui octávo Idus Aprílis pro fide cathólica 
 martyrium súbiit.
\switchcolumn
\selectlanguage{english}
\lettrine[lines=2]{S}{t.} Peter, a martyr of the Order of 
 Preachers, who was slain for the Catholic faith on the 6th day of April.
\switchcolumn*
\selectlanguage{latin}
Romæ natális sanctæ 
 Catharínæ Senénsis Vírginis, ex tértio Ordine sancti Domínici, vita et 
 miráculis claræ, quam Pius Secúndus, Póntifex Máximus, sanctárum Vírginum 
 número adscrípsit. Ipsíus tamen festum sequénti die celebrátur.
\switchcolumn
\selectlanguage{english}
At Rome, the birthday of St. 
 Catherine of Siena, virgin of the Third Order of St. Dominic, renowned for 
 her holy life and her miracles. She was inscribed among the canonized 
 virgins by Pope Pius II. Her feast, however, is celebrated on the 
 following day.
\switchcolumn*
\selectlanguage{latin}
Apud Paphum, in Cypro, 
 sancti Tychici, qui fuit discípulus beáti Pauli Apóstoli, et ab ipso 
 Apóstolo in suis Epístolis caríssimus frater, miníster fidélis suúsque in 
 Dómino consérvus appellátur.
\switchcolumn
\selectlanguage{english}
At Paphos in Cyprus, St. Tychicus, 
 a disciple of the blessed Apostle Paul, who called him in his Epistles, 
 ``most dear brother,'' ``faithful minister,'' and ``fellow-servant in the Lord'`.
\switchcolumn*
\selectlanguage{latin}
Pisis, in Túscia, 
 sancti Torpétis Mártyris, qui magnus in offício Nerónis primum fuit, unusque 
 ex his, de quibus idem Paulus Apóstolus ab urbe Roma ad Philippénses scribit: 
 « Salútant vos omnes sancti, máxime autem qui de Cǽsaris domo sunt ». 
 Sed póstea, pro fide Christi, jubénte Satéllico, 
 álapis cæditur, verbéribus 
 duríssime affícitur, ac béstiis devorándus tráditur, sed mínime læditur; 
 tandem martyrium suum decollatióne complévit.
\switchcolumn
\selectlanguage{english}
At Pisa in Tuscany, the martyr St. 
 Torpes, who filled a high office in the court of Nero, and was one of those 
 of whom the apostle wrote from Rome to the Philippians: ``All the saints 
 salute you, especially those that are of the house of Caesar.'' For the 
 faith of Christ, he was, by order of Satellicus, beaten, cruelly scourged, 
 and delivered to the beasts to be devoured, but remained uninjured. He 
 completed his martyrdom by being beheaded.
\switchcolumn*
\selectlanguage{latin}
Cirthæ, in Numídia, 
 natális sanctórum Mártyrum Agápii et Secundíni Episcopórum, qui, post longum 
 apud præfátam urbem exsílium, in persecutióne Valeriáni, in qua tunc máxime 
 Gentílium rábies ad tentándam justórum fidem inhiábat, ex illústri 
 sacerdótio effécti sunt Mártyres gloriósi. Passi sunt in eódem 
 collégio Æmiliánus miles, Tertúlla et Antónia, quæ erant sacræ 
 Vírgines, et quædam múlier cum suis géminis.
\switchcolumn
\selectlanguage{english}
At Cirta in Numidia, the birthday 
 of the holy martyrs Apapius and Secundinus, bishops, who, after a long exile 
 in that city, added to the glory of their priesthood the crown of martyrdom. 
 They suffered in the persecution of Valerian, during which the enraged 
 Gentiles made every effort to shake the faith of the just. In their 
 company suffered Aemilian, a soldier, Tertulla and Antonia, consecrated 
 virgins, and a woman with her twin children.
\switchcolumn*
\selectlanguage{latin}
In ínsula Corcyra 
 sanctórum septem Latrónum, qui, a sancto Jásone ad Christum convérsi, 
 martyrio vitam adépti sunt sempitérnam.
\switchcolumn
\selectlanguage{english}
In the island of Codyra, the seven 
 holy thieves who were converted to Christ by St. Jason, and gained eternal 
 life by martyrdom.
\switchcolumn*
\selectlanguage{latin}
Neápoli, in Campánia, 
 sancti Sevéri Epíscopi, qui, inter ália admiránda, mórtuum de sepúlcro 
 excitávit ad tempus, ut mendácem víduæ et pupíllórum creditórum argúeret 
 falsitátis.
\switchcolumn
\selectlanguage{english}
At Naples in Campania, Bishop St. 
 Severus, who, among other prodigies, raised for a short time a dead man from 
 the grave in order to convict of falsehood the lying creditor of a widow and 
 her children.
\switchcolumn*
\selectlanguage{latin}
Bríxiæ sancti Paulíni, 
 Epíscopi et Confessóris.
\switchcolumn
\selectlanguage{english}
At Brescia, St. Paulinus, bishop 
 and confessor.
\switchcolumn*
\selectlanguage{latin}
In cœnóbio Cluniacénsi, 
 in Gállia, sancti Hugónis Abbátis.
\switchcolumn
\selectlanguage{english}
In the monastery of Cluny in 
 France, St. Hugh Abbot.
\switchcolumn*
\selectlanguage{latin}
In monastério 
 Molisménsi, in Gállia, sancti Robérti, qui fuit primus Abbas Cistércii.
\switchcolumn
\selectlanguage{english}
In the monastery of Molesmes in 
 France, St. Robert, the first abbot of the Cistercians.
\switchcolumn*
\selectlanguage{latin}
\end{paracol}


% ---- martyrology/mart04/mart0430.htm
\needspace{10\baselineskip}
\begin{paracol}{2}
\selectlanguage{latin}
\begin{center}{\color{gregoriocolor} Prídie Kaléndas Maji. 
 Luna\dots\ }\end{center}
\switchcolumn
\selectlanguage{english}
\begin{center}{\color{gregoriocolor} The 
 30\textsuperscript{th} Day of April. The\dots\ Day of the Moon.}\end{center}
\end{paracol}

\noindent\begin{tabularx}{\linewidth}{*{19}{>{\centering\arraybackslash}X}}
 \textcolor{gregoriocolor}{a} & \textcolor{gregoriocolor}{b} & \textcolor{gregoriocolor}{c} & \textcolor{gregoriocolor}{d} & \textcolor{gregoriocolor}{e} & \textcolor{gregoriocolor}{f} & \textcolor{gregoriocolor}{g} & \textcolor{gregoriocolor}{h} & \textcolor{gregoriocolor}{i} & \textcolor{gregoriocolor}{k} & \textcolor{gregoriocolor}{l} & \textcolor{gregoriocolor}{m} & \textcolor{gregoriocolor}{n} & \textcolor{gregoriocolor}{p} & \textcolor{gregoriocolor}{q} & \textcolor{gregoriocolor}{r} & \textcolor{gregoriocolor}{s} & \textcolor{gregoriocolor}{t} & \textcolor{gregoriocolor}{u} \\
 3 & 4 & 5 & 6 & 7 & 8 & 9 & 10 & 11 & 12 & 13 & 14 & 15 & 16 & 17 & 18 & 19 & 20 & 21 \\
\end{tabularx}
\vspace{0.5\baselineskip}
\noindent\begin{tabularx}{\linewidth}{*{12}{>{\centering\arraybackslash}X}}
 \textcolor{gregoriocolor}{A} & \textcolor{gregoriocolor}{B} & \textcolor{gregoriocolor}{C} & \textcolor{gregoriocolor}{D} & \textcolor{gregoriocolor}{E} & F & \textcolor{gregoriocolor}{F} & \textcolor{gregoriocolor}{G} & \textcolor{gregoriocolor}{H} & \textcolor{gregoriocolor}{M} & \textcolor{gregoriocolor}{N} & \textcolor{gregoriocolor}{P} \\
 22 & 23 & 24 & 25 & 26 & 27 & 26 & 27 & 28 & 29 & 1 & 2 \\
\end{tabularx}

\begin{paracol}{2}
\selectlanguage{latin}
\lettrine[lines=2]{S}{anctæ} Catharínæ 
 Senénsis Vírginis, ex tértio Ordine sancti Domínici, quæ ad cæléstem Sponsum 
 transívit prídie hujus diéi.
\switchcolumn
\selectlanguage{english}
\lettrine[lines=2]{S}{t.} Catherine of Siena, virgin of 
 the Third Order of St. Dominic, who on the previous day went to her heavenly 
 Spouse.
\switchcolumn*
\selectlanguage{latin}
Apud Sántonas, in 
 Gállia, beáti Eutrópii, Epíscopi et Mártyris, qui, a sancto Cleménte Papa, 
 órdinis pontificális grátia consecrátus, in Gálliam diréctus est, ibíque, 
 perácta diu prædicatióne, tandem, ob Christi testimónium, collíso cápite, 
 victor occúbuit.
\switchcolumn
\selectlanguage{english}
At Saintes in France, blessed 
 Eutropius, bishop and martyr, who was consecrated bishop and sent to France 
 by St. Clement. After preaching for many years, he had his skull 
 crushed for bearing testimony to Christ, and thus gained victory by his 
 death.
\switchcolumn*
\selectlanguage{latin}
Córdubæ, in Hispánia, sanctórum Mártyrum Amatóris Presbyteri, Petri Mónachi, et Ludovíci.
\switchcolumn
\selectlanguage{english}
At Cordova in Spain, the holy 
 martyrs Amator, a priest, Peter, a monk, and Louis.
\switchcolumn*
\selectlanguage{latin}
Nováriæ sancti 
 Lauréntii Presbyteri, et puerórum Mártyrum, quos ille suscéperat educándos.
\switchcolumn
\selectlanguage{english}
At Novara, the martyrdom of the 
 holy priest Laurence, and some boys whom he was teaching.
\switchcolumn*
\selectlanguage{latin}
Alexandríæ sanctórum 
 Mártyrum Aphrodísii Presbyteri, et aliórum trigínta.
\switchcolumn
\selectlanguage{english}
At Alexandria, the holy martyrs 
 Aphrodisius, a priest, and thirty martyrs.
\switchcolumn*
\selectlanguage{latin}
Lambésæ, in Numídia, 
 natális sanctórum Mártyrum Mariáni Lectóris, et Jacóbi Diáconi. Horum 
 prior, cum infestatiónes jamprídem Deciánæ persecutiónis in confessióne 
 Christi vicísset, íterum cum claríssimo colléga tentus est; et ambo, post 
 dira et exquisíta supplícia, mirabíliter divínis revelatiónibus secúndo 
 confortáti, novíssime, cum multis áliis, gládio consummáti sunt.
\switchcolumn
\selectlanguage{english}
At Lambesa in Numidia, the birthday 
 of the holy martyrs Marian, a lector, and James, a deacon. The former, 
 after having successfully endured many trials for the confession of Christ 
 in the persecution of Decius, was again arrested with his noble companions, 
 and both were subjected to severe and cruel torments, during which they were 
 twice miraculously comforted by heaven, but finally fell by the sword along 
 with many others.
\switchcolumn*
\selectlanguage{latin}
Ephesi sancti Máximi 
 Mártyris, qui in persecutióne Décii coronátus est.
\switchcolumn
\selectlanguage{english}
At Ephesus, the martyr St. Maximus, 
 who received his crown during the persecution of Decius.
\switchcolumn*
\selectlanguage{latin}
Firmi, in Picéno, 
 sanctæ Sophíæ, Vírginis et Mártyris.
\switchcolumn
\selectlanguage{english}
At Ferno in Piceno, St. Sophia, 
 virgin and martyr.
\switchcolumn*
\selectlanguage{latin}
Evóreæ, in Epíro, 
 sancti Donáti Epíscopi, qui, témpore Theodósii Imperatóris, exímia 
 sanctitáte refúlsit.
\switchcolumn
\selectlanguage{english}
At Evorea in Epirus, St. Donatus, a 
 bishop, who was eminent for sanctity in the time of Emperor Theodosius.
\switchcolumn*
\selectlanguage{latin}
Neápoli, in Campánia, sancti Pompónii Epíscopi.
\switchcolumn
\selectlanguage{english}
At Naples in Campania, St. 
 Pomponius, bishop.
\switchcolumn*
\selectlanguage{latin}
Londíni, in Anglia, 
 sancti Erconváldi Epíscopi, qui multis miráculis cláruit.
\switchcolumn
\selectlanguage{english}
At London in England, St. Erkenwald, 
 a bishop celebrated for many miracles.
\switchcolumn*
\selectlanguage{latin}
Chérii, apud Augústam 
 Taurinórum, sancti Joséphi Benedícti Cottoléngo, Confessóris, Parvæ Domus a 
 Divína Providéntia Fundatóris, summa in Deum confidéntia et caritáte in 
 páuperes insígnis, quem Pius Papa Undécimus Sanctórum fastis adscrípsit.
\switchcolumn
\selectlanguage{english}
At Chieri, near Turin, St. Joseph 
 Cottolengo, confessor, founder of the Little House of Divine Providence, 
 full of trust in God and remarkable for his charity toward the poor, whom 
 Pope Pius XI enrolled among the saints.
\switchcolumn*
\selectlanguage{latin}
\end{paracol}
